% =============================================================================
%  progress-report-05.tex — รายงานความก้าวหน้าโครงการ ULMs ครั้งที่ 5
%  *** ฉบับส่งอาจารย์ รอบ 18 – 24 ก.ย. 2569 (วันที่ 14 จาก 21 ของ Sprint 4) ***
%
%  build:  cd ~/Projects/SoftwareEn/docs && latexmk progress-report-05.tex
%          ผลลัพธ์ออกที่ build/progress-report-05.pdf
%
%  หลักการเขียน — ยกมาจาก progress-report-03.tex และ -04.tex ทุกข้อ
%    สั้น ไม่มีสารบัญ · แผนเส้นเดียว · อ้างได้เฉพาะฉบับที่ส่งจริง (ครั้งที่ 1 ถึง 4)
%    ห้ามรายงานเป็นจำนวนบรรทัด · ไม่มี \textbf กลางย่อหน้า
%    ศัพท์ของการพัฒนาใช้อังกฤษ และเว้นวรรคหน้าหลังทุกครั้งที่ชนอักษรไทย
%    ห้ามรัน regex ยุบช่องไฟทับทั้งไฟล์
%
%  ห้ามแก้ progress-report-01/02/03/04.tex ย้อนหลัง ทั้งสี่ฉบับส่งไปแล้ว
%
%  ตัวเลขทุกตัวอ้าง origin/main ที่ e0a4865 (24 ก.ย. 2569) fetch ซ้ำก่อนสร้าง PDF แล้ว
%  ไม่นับ spike-option-b/ และไม่นับ branch ที่ยังไม่ merge
%
%  *** กำหนดเสร็จเลื่อนออกไป 1 สัปดาห์ เป็น 8 ต.ค. 2569 ***
%  ผู้รับผิดชอบโครงการอนุมัติให้ยืด Sprint 4 ออกไปอีกหนึ่งสัปดาห์ ไม่ใช่ยืด Sprint 5
%    Sprint 4 = 11 ก.ย. – 1 ต.ค. (21 วัน) · Sprint 5 = 2 – 8 ต.ค. (7 วัน)
%  วันรายงานจึงไม่ใช่ขอบ Sprint อีกต่อไป แต่เป็นวันที่ 14 จาก 21 ของ Sprint 4
%
%  *** ห้ามเลื่อนเส้นแผนตามกำหนดเสร็จ ***
%  ผู้รับผิดชอบโครงการกำหนดว่าสัปดาห์ที่เพิ่มมาเป็นเวลาไล่ตามงานที่ยังไม่เสร็จ
%  ไม่ใช่การเลื่อนเป้าหมาย ค่าเป้าหมายรายวันในภาคผนวกจึงคงเดิมทุกช่อง
%  แผน ณ 24 ก.ย. ยังเป็น 94% และแผนยังแตะ 100% ที่ 1 ต.ค. เหมือนเดิม
%  ช่วง 2 – 8 ต.ค. เส้นแผนราบอยู่ที่ 100 เพราะไม่มีเป้าใหม่ มีแต่เวลาให้ไล่ตาม
%  อย่าสอดแทรกค่าแผนใหม่ให้วันรายงาน เคยคิดแบบนั้นแล้วได้ 89% ซึ่งทำให้
%  ระยะห่างดูแคบลงทั้งที่ทีมไม่ได้ทำงานได้มากขึ้น — ผู้ใช้สั่งไม่เอาวิธีนั้น
%
%  *** ดึงงานที่ตัดไปกลับเข้าขอบเขต ***
%  ตะกร้าข้ามหน่วยงานและการจองล่วงหน้า ซึ่งรายงานครั้งที่ 4 ตัดออกตามแผนสำรอง
%  ถูกดึงกลับเข้า Sprint 4 ที่ยืดออก งานค้างจึงโตขึ้น แต่สองรายการนี้สถานะไม่เท่ากัน
%    ตะกร้าข้ามหน่วยงาน — ไม่มีโค้ดเลยทั้งสองฝั่ง ร่างคำขอไม่จัดกลุ่มตามหน่วยงาน
%                        และ loan.request.service.ts ไม่มีตรรกะแยกหน่วยงาน
%    การจองล่วงหน้า    — หน้าเว็บทำแล้ว ช่องเลือกวันรับของจำกัดที่ RESERVATION_MAX_DAYS
%                        (90 วัน) แต่เซิร์ฟเวอร์ไม่บังคับ RESERVATION_TOO_FAR มีแต่คำแปลไทย
%  BE Dev จึงลดจาก 93 เป็น 90 (ขาดทั้งสองรายการ) ส่วน FE Dev ลดจาก 92 เป็น 90
%  (ขาดเฉพาะตะกร้าข้ามหน่วยงาน) อย่าเขียนว่า "ทั้งสองยังไม่มีโค้ด" — ไม่จริง
%
%  *** วิธีนับเปลี่ยนที่ฉบับนี้ ***
%  ฉบับ 1 ถึง 4 คิดความก้าวหน้าจากการนับของที่สร้างเสร็จ (procedure หน้าจอ ชุดทดสอบ)
%  ฉบับนี้หักข้อบกพร่องที่พบจากการตรวจ UI ด้วยมือเมื่อ 19 ก.ย. ซึ่งเกิดขึ้นหลัง
%  รายงานครั้งที่ 4 (17 ก.ย.) จึงไม่เคยถูกหักมาก่อน ผลคือ 86% แทน 90% ที่ได้จากการนับล้วน
%  ไม่ย้อนแก้ค่าของฉบับที่ส่งไปแล้ว จุดบนเส้นผลจริงก่อน 24 ก.ย. จึงยังเป็นค่าเดิม
%  บัญชีข้อบกพร่องทั้ง 30 ข้ออยู่ในหัวข้อ 7
%
%  ที่มาของตัวเลข ตรวจซ้ำได้จาก repo:
%    102 procedure / 11 domain = ผลรวม @Query/@Mutation ใน backend/src/*/*.router.ts
%                    ยืนยันซ้ำด้วย `npm run trpc:generate` ซึ่งรายงาน 11 router 102 procedure
%                    (รอบก่อน 86 / 10 — domain ที่เพิ่มคือ appeal)
%    33 ตาราง 9 enum 15 migration = backend/prisma/ (รอบก่อน 30 / 9 / 10)
%    421 test / 41 suite = `npm test` ในเครื่องนี้ ชี้ DATABASE_URL ไปฐาน ulms_test
%                    ผ่านทั้งหมด (รอบก่อน 266 / 24)
%                    หมายเหตุ ต้อง `npm install` ก่อน เพราะรอบนี้เพิ่ม nodemailer
%    19 จาก 27 service มี .spec.ts คู่กัน (รอบก่อน 10 จาก 23 วิธีนับเดียวกัน)
%    105 vitest / 17 ไฟล์ (รอบก่อน 33 / 3) · 49 AC check เท่าเดิม
%    lint ฝั่ง frontend 5 error 8 warning (รอบก่อน 0 error 7 warning)
%                    error ทั้ง 5 อยู่ใน frontend/tests/auth/auth-flow.test.tsx
%                    CI ไม่มีขั้น lint จึงไม่จับ
%    20 ผ่าน 5 ไม่ผ่าน จาก 25 = `npx playwright test` โดยสตาร์ต backend :3000 และ
%                    frontend :5173 เองบนฐาน ulms_e2e ที่ seed ด้วย `npm run seed`
%                    (ไม่ใช่ `npm run db:seed` ซึ่งเป็นคนละชุดข้อมูลและไม่มี test_borrower)
%    34 หน้า = frontend/src/features/**/*-page.tsx (รอบก่อน 29)
%    34 มี UI จริง = ไม่มีไฟล์ไหนใช้ PlaceholderPage แล้ว (รอบก่อนเหลือหน้าว่าง 2 หน้า)
%    0 หน้าที่แสดงข้อมูลจำลอง = CATALOG_ITEMS / AUDIT_EVENTS / MOCK_APPEALS ไม่ถูกอ้างถึง
%                    จากที่ใดนอก mock-data.ts แล้ว ที่เหลือใน mock-data.ts เป็น type
%                    ค่าคงที่ร่วม และฟังก์ชันช่วย
%    102 procedure / 11 domain ที่ FE ประกาศใน frontend/src/server/trpc-contract.ts
%                    ตรงกับเซิร์ฟเวอร์ทุกชื่อ ต่างกัน 0 รายการ (รอบก่อนต่างกัน 14)
%    82 จาก 89 รหัสข้อผิดพลาดมีคำแปลไทย (รอบก่อน 62 จาก 65)
%                    7 ตัวที่ขาดเป็นของ appeal และ inspection ที่เพิ่งเพิ่มรอบนี้
%    6 cron job ทำงานจริงทั้ง 6 (รอบก่อน 5 จาก 8 — computeAvailability กับ
%                    rollupDailyStats ถูกถอดออกโดยตั้งใจใน 6ac43ce)
%    71 commit / 23 PR ในช่วง 0cb21a7..e0a4865 (ไม่รวม merge = 45)
%    ผลตรวจ UI 19 ก.ย. = เด็ค Canva 20 หน้า 33 ข้อ (30 ข้อบกพร่อง + 3 คำถาม)
%                    แก้แล้ว 16 (ทีมกำกับเอง 8 + ตรวจยืนยันจากโค้ดอีก 8)
%                    แก้บางส่วน 5 · ยังไม่แก้ 9
%                    ข้อที่ยืนยันชัดที่สุดว่ายังไม่แก้คือการเรียงยอดนิยม ซึ่ง
%                    catalog-page.tsx ยังเรียงด้วย b.totalUnits - a.totalUnits
% =============================================================================
% =============================================================================
%  ulms-preamble.tex — preamble ที่ใช้ร่วมกันทุกเอกสารของโครงการ ULMs
%
%  ใช้โดย:  main.tex (ข้อเสนอโครงการ) · progress.tex (รายงานความก้าวหน้า)
%
%  เอกสารที่เรียกใช้ต้องทำสองอย่างหลัง % =============================================================================
%  ulms-preamble.tex — preamble ที่ใช้ร่วมกันทุกเอกสารของโครงการ ULMs
%
%  ใช้โดย:  main.tex (ข้อเสนอโครงการ) · progress.tex (รายงานความก้าวหน้า)
%
%  เอกสารที่เรียกใช้ต้องทำสองอย่างหลัง % =============================================================================
%  ulms-preamble.tex — preamble ที่ใช้ร่วมกันทุกเอกสารของโครงการ ULMs
%
%  ใช้โดย:  main.tex (ข้อเสนอโครงการ) · progress.tex (รายงานความก้าวหน้า)
%
%  เอกสารที่เรียกใช้ต้องทำสองอย่างหลัง \input{ulms-preamble}:
%    1) \hypersetup{pdftitle={...}}  — ชื่อเอกสารใน metadata ของ PDF
%    2) \begin{document}
%  (pdftitle ถูกย้ายออกจากไฟล์นี้ เพราะเป็นค่าเฉพาะของแต่ละเอกสาร)
%
%  แก้ที่นี่ที่เดียว = มีผลกับทุกเอกสาร ห้ามก๊อป preamble ไปวางซ้ำ
% =============================================================================
% =============================================================================
%  main.tex  —  เอกสารข้อเสนอโครงการ (Academic Project Proposal)
%  ระบบบริหารจัดการการยืม–คืนอุปกรณ์มหาวิทยาลัย (ULMs)
%
%  IMPORTANT: This document contains Thai text and MUST be compiled with
%  XeLaTeX (not pdflatex). Thai needs a Unicode font + ICU line-breaking.
%      Build:  latexmk main.tex        (uses .latexmkrc -> xelatex)
%      or:     xelatex main.tex        (run twice for the table of contents)
% =============================================================================

\documentclass[11pt, a4paper]{article}

% ---- Fonts & Thai support ----------------------------------------------------
% fontspec lets XeLaTeX use any system font. Laksaman is a complete serif face
% that covers BOTH Thai and Latin glyphs, so one \setmainfont handles the whole
% (Thai + English) document without font switching.
\usepackage{fontspec}

% Thai has no spaces between words, so TeX cannot find line-break points on its
% own. These two XeTeX primitives switch on ICU's dictionary-based Thai word
% segmentation, which is what makes Thai paragraphs wrap correctly.
\XeTeXlinebreaklocale "th"
\XeTeXlinebreakskip = 0pt plus 0.1pt

\setmainfont{Laksaman}

% ---- Page geometry & paragraph style ----------------------------------------
\usepackage[margin=2.5cm]{geometry}
\usepackage{parskip}          % block paragraphs (no indent, gap between) — common
                              % and more readable for Thai documents
\linespread{1.15}             % a little extra leading; Thai has tall tone marks

% ---- Colours (monochrome / grayscale palette) ------------------------------
\usepackage[table]{xcolor}    % 'table' option enables \rowcolor for table headers
% Monochrome (grayscale) palette. Names kept as-is so every reference still
% works; only the tones changed from the original green/blue accents.
\definecolor{kugreen}{HTML}{1A1A1A}   % near-black — section headings & title
\definecolor{kudark}{HTML}{333333}    % dark gray — subsection/subsubsection titles
\definecolor{headbg}{HTML}{EAEAEA}    % light gray background for table header rows
\definecolor{linkblue}{HTML}{404040}  % medium-dark gray — links / citations

% ---- Section heading styling -------------------------------------------------
% 'nobottomtitles*' stops a heading from being printed near the bottom of a
% page: if less than \bottomtitlespace of room is left, the heading (and its
% following text) moves to the next page. This is titlesec's built-in, robust
% replacement for the earlier \needspace hack — the latter injected vertical
% glue that clashed with longtable page-splitting ("Infinite glue shrinkage").
\usepackage[nobottomtitles*]{titlesec}
\titleformat{\section}
  {\color{kugreen}\Large\bfseries}{\thesection}{0.6em}{}
\titleformat{\subsection}
  {\color{kudark}\large\bfseries}{\thesubsection}{0.5em}{}
\titleformat{\subsubsection}
  {\color{kudark}\normalsize\bfseries}{\thesubsubsection}{0.5em}{}
\titlespacing*{\section}{0pt}{1.4\baselineskip}{0.5\baselineskip}

% Minimum room a heading needs at the bottom of a page; if only ~1–3 lines are
% left (less than this), the heading is bumped to the next page instead of being
% stranded above its table/text.
\renewcommand{\bottomtitlespace}{4\baselineskip}

% ---- Hierarchical indentation by heading level ------------------------------
% Each heading level shifts BOTH the heading and the body text under it further
% to the right: section = 0, subsection (x.y) = 1 step, subsubsection (x.y.z) =
% 2 steps. We do this by having the sectioning commands set \leftskip (a global
% left margin that carries through every following line until the next heading
% resets it). \par first, so the *previous* paragraph is not re-indented.
% \leftskip indents plain paragraphs; \@totalleftmargin makes list environments
% (itemize/enumerate) start from the same indent instead of the page margin.
\makeatletter
\newlength{\lvlindent}
\setlength{\lvlindent}{1.8em}          % one indent step
\let\ulmsSecIndent\section
\renewcommand{\section}{\par\global\leftskip=0pt\global\@totalleftmargin=0pt\relax\ulmsSecIndent}
\let\ulmsSubsecIndent\subsection
\renewcommand{\subsection}{\par\global\leftskip=\lvlindent\global\@totalleftmargin=\lvlindent\relax\ulmsSubsecIndent}
\let\ulmsSubsubsecIndent\subsubsection
\renewcommand{\subsubsection}{\par\global\leftskip=2\lvlindent\global\@totalleftmargin=2\lvlindent\relax\ulmsSubsubsecIndent}
\makeatother

% Indent the headings themselves to match (titlesec ignores \leftskip for the
% heading line, so its left margin is set here via the first \titlespacing arg).
\titlespacing*{\subsection}   {\lvlindent} {1.0\baselineskip}{0.4\baselineskip}
\titlespacing*{\subsubsection}{2\lvlindent}{0.8\baselineskip}{0.3\baselineskip}

% ---- Lists (tighter, custom bullets) ----------------------------------------
\usepackage{enumitem}
\setlist[itemize]{leftmargin=1.6em, itemsep=2pt, topsep=3pt}
\setlist[enumerate]{leftmargin=1.9em, itemsep=2pt, topsep=3pt}

% ---- Tables ------------------------------------------------------------------
\usepackage{array}
\usepackage{booktabs}         % \toprule / \midrule / \bottomrule
\usepackage{tabularx}         % auto-width columns that wrap long Thai text
\usepackage{longtable}        % tables that can break across pages
\usepackage{multirow}         % merged cells for the team-structure table

% Left-aligned wrapping column types:
%   L{width} — fixed-width ragged-right paragraph column
%   Y        — flexible ragged-right X column (for tabularx)
\newcolumntype{L}[1]{>{\raggedright\arraybackslash}p{#1}}
\newcolumntype{Y}{>{\raggedright\arraybackslash}X}
\renewcommand{\arraystretch}{1.3}

% Helpers for a styled header row inside tables.
%   \hrow  — starts a header row and shades it (must be FIRST in the row;
%            colortbl allows only one \rowcolor per row).
%   \thead — bold header-cell text.
\newcommand{\hrow}{\rowcolor{headbg}}
\newcommand{\thead}[1]{\textbf{#1}}

% \begin{keeptable} … \end{keeptable} — keeps a table and the bold label above
% it on the SAME page (used by the data dictionary in §5). A tabularx cannot
% split across pages, so LaTeX is free to break between the label and its
% table, stranding the label alone at the bottom of a page; wrapping the pair
% in a minipage makes them one unbreakable block.
%   * \leftskip is zeroed inside: the minipage box is already placed at the
%     current heading indent, so leaving it set would indent the content twice
%     and push the table past the right margin.
%   * \parskip is restored because minipage's \@parboxrestore zeroes it, which
%     would glue the table straight onto the label line.
% Tables inside must use \linewidth (= the minipage width), not \textwidth.
\newenvironment{keeptable}
  {\par\medskip\noindent
   \begin{minipage}{\dimexpr\textwidth-\leftskip\relax}%
   \leftskip=0pt\relax
   \setlength{\parskip}{0.5\baselineskip plus 2pt}}
  {\end{minipage}\par\medskip}

% ---- Table of contents styling ----------------------------------------------
% By default LaTeX gives dotted leaders to subsections but NOT to sections
% (they print bold with just a right-aligned page number). tocloft lets us turn
% on dots for the section level too, so every TOC line runs "….." to its page.
\usepackage{tocloft}
\renewcommand{\cftsecdotsep}{\cftdotsep}   % dotted leaders for section entries
\renewcommand{\cftsecleader}{\cftdotfill{\cftsecdotsep}}
\renewcommand{\cftsecpagefont}{\normalfont}  % keep page numbers un-bold

% ---- Figures -----------------------------------------------------------------
\usepackage{graphicx}
\graphicspath{{figures/}}

% Landscape (rotated) full-page floats — needed by the ER diagram in §5, which
% is far wider than it is tall and would be unreadable scaled to \textwidth.
\usepackage{rotating}

% ---- Flowcharts / diagrams (TikZ) -------------------------------------------
% Used for the borrow–return workflow diagram (§5). Libraries: geometric shapes
% for decision diamonds, arrows.meta for arrowheads, positioning for relative
% node placement, calc/fit/backgrounds for routing loops and grouping.
\usepackage{tikz}
\usetikzlibrary{shapes.geometric, arrows.meta, positioning, calc, fit, backgrounds}

% ---- Links & bookmarks -------------------------------------------------------
\usepackage{hyperref}
\hypersetup{
  colorlinks=true, linkcolor=black, citecolor=linkblue, urlcolor=linkblue,
  pdflang={th},
}

% Thai labels for auto-generated headings.
\renewcommand{\contentsname}{สารบัญ (Table of Contents)}
\renewcommand{\tablename}{ตาราง}
\renewcommand{\figurename}{รูปที่}

:
%    1) \hypersetup{pdftitle={...}}  — ชื่อเอกสารใน metadata ของ PDF
%    2) \begin{document}
%  (pdftitle ถูกย้ายออกจากไฟล์นี้ เพราะเป็นค่าเฉพาะของแต่ละเอกสาร)
%
%  แก้ที่นี่ที่เดียว = มีผลกับทุกเอกสาร ห้ามก๊อป preamble ไปวางซ้ำ
% =============================================================================
% =============================================================================
%  main.tex  —  เอกสารข้อเสนอโครงการ (Academic Project Proposal)
%  ระบบบริหารจัดการการยืม–คืนอุปกรณ์มหาวิทยาลัย (ULMs)
%
%  IMPORTANT: This document contains Thai text and MUST be compiled with
%  XeLaTeX (not pdflatex). Thai needs a Unicode font + ICU line-breaking.
%      Build:  latexmk main.tex        (uses .latexmkrc -> xelatex)
%      or:     xelatex main.tex        (run twice for the table of contents)
% =============================================================================

\documentclass[11pt, a4paper]{article}

% ---- Fonts & Thai support ----------------------------------------------------
% fontspec lets XeLaTeX use any system font. Laksaman is a complete serif face
% that covers BOTH Thai and Latin glyphs, so one \setmainfont handles the whole
% (Thai + English) document without font switching.
\usepackage{fontspec}

% Thai has no spaces between words, so TeX cannot find line-break points on its
% own. These two XeTeX primitives switch on ICU's dictionary-based Thai word
% segmentation, which is what makes Thai paragraphs wrap correctly.
\XeTeXlinebreaklocale "th"
\XeTeXlinebreakskip = 0pt plus 0.1pt

\setmainfont{Laksaman}

% ---- Page geometry & paragraph style ----------------------------------------
\usepackage[margin=2.5cm]{geometry}
\usepackage{parskip}          % block paragraphs (no indent, gap between) — common
                              % and more readable for Thai documents
\linespread{1.15}             % a little extra leading; Thai has tall tone marks

% ---- Colours (monochrome / grayscale palette) ------------------------------
\usepackage[table]{xcolor}    % 'table' option enables \rowcolor for table headers
% Monochrome (grayscale) palette. Names kept as-is so every reference still
% works; only the tones changed from the original green/blue accents.
\definecolor{kugreen}{HTML}{1A1A1A}   % near-black — section headings & title
\definecolor{kudark}{HTML}{333333}    % dark gray — subsection/subsubsection titles
\definecolor{headbg}{HTML}{EAEAEA}    % light gray background for table header rows
\definecolor{linkblue}{HTML}{404040}  % medium-dark gray — links / citations

% ---- Section heading styling -------------------------------------------------
% 'nobottomtitles*' stops a heading from being printed near the bottom of a
% page: if less than \bottomtitlespace of room is left, the heading (and its
% following text) moves to the next page. This is titlesec's built-in, robust
% replacement for the earlier \needspace hack — the latter injected vertical
% glue that clashed with longtable page-splitting ("Infinite glue shrinkage").
\usepackage[nobottomtitles*]{titlesec}
\titleformat{\section}
  {\color{kugreen}\Large\bfseries}{\thesection}{0.6em}{}
\titleformat{\subsection}
  {\color{kudark}\large\bfseries}{\thesubsection}{0.5em}{}
\titleformat{\subsubsection}
  {\color{kudark}\normalsize\bfseries}{\thesubsubsection}{0.5em}{}
\titlespacing*{\section}{0pt}{1.4\baselineskip}{0.5\baselineskip}

% Minimum room a heading needs at the bottom of a page; if only ~1–3 lines are
% left (less than this), the heading is bumped to the next page instead of being
% stranded above its table/text.
\renewcommand{\bottomtitlespace}{4\baselineskip}

% ---- Hierarchical indentation by heading level ------------------------------
% Each heading level shifts BOTH the heading and the body text under it further
% to the right: section = 0, subsection (x.y) = 1 step, subsubsection (x.y.z) =
% 2 steps. We do this by having the sectioning commands set \leftskip (a global
% left margin that carries through every following line until the next heading
% resets it). \par first, so the *previous* paragraph is not re-indented.
% \leftskip indents plain paragraphs; \@totalleftmargin makes list environments
% (itemize/enumerate) start from the same indent instead of the page margin.
\makeatletter
\newlength{\lvlindent}
\setlength{\lvlindent}{1.8em}          % one indent step
\let\ulmsSecIndent\section
\renewcommand{\section}{\par\global\leftskip=0pt\global\@totalleftmargin=0pt\relax\ulmsSecIndent}
\let\ulmsSubsecIndent\subsection
\renewcommand{\subsection}{\par\global\leftskip=\lvlindent\global\@totalleftmargin=\lvlindent\relax\ulmsSubsecIndent}
\let\ulmsSubsubsecIndent\subsubsection
\renewcommand{\subsubsection}{\par\global\leftskip=2\lvlindent\global\@totalleftmargin=2\lvlindent\relax\ulmsSubsubsecIndent}
\makeatother

% Indent the headings themselves to match (titlesec ignores \leftskip for the
% heading line, so its left margin is set here via the first \titlespacing arg).
\titlespacing*{\subsection}   {\lvlindent} {1.0\baselineskip}{0.4\baselineskip}
\titlespacing*{\subsubsection}{2\lvlindent}{0.8\baselineskip}{0.3\baselineskip}

% ---- Lists (tighter, custom bullets) ----------------------------------------
\usepackage{enumitem}
\setlist[itemize]{leftmargin=1.6em, itemsep=2pt, topsep=3pt}
\setlist[enumerate]{leftmargin=1.9em, itemsep=2pt, topsep=3pt}

% ---- Tables ------------------------------------------------------------------
\usepackage{array}
\usepackage{booktabs}         % \toprule / \midrule / \bottomrule
\usepackage{tabularx}         % auto-width columns that wrap long Thai text
\usepackage{longtable}        % tables that can break across pages
\usepackage{multirow}         % merged cells for the team-structure table

% Left-aligned wrapping column types:
%   L{width} — fixed-width ragged-right paragraph column
%   Y        — flexible ragged-right X column (for tabularx)
\newcolumntype{L}[1]{>{\raggedright\arraybackslash}p{#1}}
\newcolumntype{Y}{>{\raggedright\arraybackslash}X}
\renewcommand{\arraystretch}{1.3}

% Helpers for a styled header row inside tables.
%   \hrow  — starts a header row and shades it (must be FIRST in the row;
%            colortbl allows only one \rowcolor per row).
%   \thead — bold header-cell text.
\newcommand{\hrow}{\rowcolor{headbg}}
\newcommand{\thead}[1]{\textbf{#1}}

% \begin{keeptable} … \end{keeptable} — keeps a table and the bold label above
% it on the SAME page (used by the data dictionary in §5). A tabularx cannot
% split across pages, so LaTeX is free to break between the label and its
% table, stranding the label alone at the bottom of a page; wrapping the pair
% in a minipage makes them one unbreakable block.
%   * \leftskip is zeroed inside: the minipage box is already placed at the
%     current heading indent, so leaving it set would indent the content twice
%     and push the table past the right margin.
%   * \parskip is restored because minipage's \@parboxrestore zeroes it, which
%     would glue the table straight onto the label line.
% Tables inside must use \linewidth (= the minipage width), not \textwidth.
\newenvironment{keeptable}
  {\par\medskip\noindent
   \begin{minipage}{\dimexpr\textwidth-\leftskip\relax}%
   \leftskip=0pt\relax
   \setlength{\parskip}{0.5\baselineskip plus 2pt}}
  {\end{minipage}\par\medskip}

% ---- Table of contents styling ----------------------------------------------
% By default LaTeX gives dotted leaders to subsections but NOT to sections
% (they print bold with just a right-aligned page number). tocloft lets us turn
% on dots for the section level too, so every TOC line runs "….." to its page.
\usepackage{tocloft}
\renewcommand{\cftsecdotsep}{\cftdotsep}   % dotted leaders for section entries
\renewcommand{\cftsecleader}{\cftdotfill{\cftsecdotsep}}
\renewcommand{\cftsecpagefont}{\normalfont}  % keep page numbers un-bold

% ---- Figures -----------------------------------------------------------------
\usepackage{graphicx}
\graphicspath{{figures/}}

% Landscape (rotated) full-page floats — needed by the ER diagram in §5, which
% is far wider than it is tall and would be unreadable scaled to \textwidth.
\usepackage{rotating}

% ---- Flowcharts / diagrams (TikZ) -------------------------------------------
% Used for the borrow–return workflow diagram (§5). Libraries: geometric shapes
% for decision diamonds, arrows.meta for arrowheads, positioning for relative
% node placement, calc/fit/backgrounds for routing loops and grouping.
\usepackage{tikz}
\usetikzlibrary{shapes.geometric, arrows.meta, positioning, calc, fit, backgrounds}

% ---- Links & bookmarks -------------------------------------------------------
\usepackage{hyperref}
\hypersetup{
  colorlinks=true, linkcolor=black, citecolor=linkblue, urlcolor=linkblue,
  pdflang={th},
}

% Thai labels for auto-generated headings.
\renewcommand{\contentsname}{สารบัญ (Table of Contents)}
\renewcommand{\tablename}{ตาราง}
\renewcommand{\figurename}{รูปที่}

:
%    1) \hypersetup{pdftitle={...}}  — ชื่อเอกสารใน metadata ของ PDF
%    2) \begin{document}
%  (pdftitle ถูกย้ายออกจากไฟล์นี้ เพราะเป็นค่าเฉพาะของแต่ละเอกสาร)
%
%  แก้ที่นี่ที่เดียว = มีผลกับทุกเอกสาร ห้ามก๊อป preamble ไปวางซ้ำ
% =============================================================================
% =============================================================================
%  main.tex  —  เอกสารข้อเสนอโครงการ (Academic Project Proposal)
%  ระบบบริหารจัดการการยืม–คืนอุปกรณ์มหาวิทยาลัย (ULMs)
%
%  IMPORTANT: This document contains Thai text and MUST be compiled with
%  XeLaTeX (not pdflatex). Thai needs a Unicode font + ICU line-breaking.
%      Build:  latexmk main.tex        (uses .latexmkrc -> xelatex)
%      or:     xelatex main.tex        (run twice for the table of contents)
% =============================================================================

\documentclass[11pt, a4paper]{article}

% ---- Fonts & Thai support ----------------------------------------------------
% fontspec lets XeLaTeX use any system font. Laksaman is a complete serif face
% that covers BOTH Thai and Latin glyphs, so one \setmainfont handles the whole
% (Thai + English) document without font switching.
\usepackage{fontspec}

% Thai has no spaces between words, so TeX cannot find line-break points on its
% own. These two XeTeX primitives switch on ICU's dictionary-based Thai word
% segmentation, which is what makes Thai paragraphs wrap correctly.
\XeTeXlinebreaklocale "th"
\XeTeXlinebreakskip = 0pt plus 0.1pt

\setmainfont{Laksaman}

% ---- Page geometry & paragraph style ----------------------------------------
\usepackage[margin=2.5cm]{geometry}
\usepackage{parskip}          % block paragraphs (no indent, gap between) — common
                              % and more readable for Thai documents
\linespread{1.15}             % a little extra leading; Thai has tall tone marks

% ---- Colours (monochrome / grayscale palette) ------------------------------
\usepackage[table]{xcolor}    % 'table' option enables \rowcolor for table headers
% Monochrome (grayscale) palette. Names kept as-is so every reference still
% works; only the tones changed from the original green/blue accents.
\definecolor{kugreen}{HTML}{1A1A1A}   % near-black — section headings & title
\definecolor{kudark}{HTML}{333333}    % dark gray — subsection/subsubsection titles
\definecolor{headbg}{HTML}{EAEAEA}    % light gray background for table header rows
\definecolor{linkblue}{HTML}{404040}  % medium-dark gray — links / citations

% ---- Section heading styling -------------------------------------------------
% 'nobottomtitles*' stops a heading from being printed near the bottom of a
% page: if less than \bottomtitlespace of room is left, the heading (and its
% following text) moves to the next page. This is titlesec's built-in, robust
% replacement for the earlier \needspace hack — the latter injected vertical
% glue that clashed with longtable page-splitting ("Infinite glue shrinkage").
\usepackage[nobottomtitles*]{titlesec}
\titleformat{\section}
  {\color{kugreen}\Large\bfseries}{\thesection}{0.6em}{}
\titleformat{\subsection}
  {\color{kudark}\large\bfseries}{\thesubsection}{0.5em}{}
\titleformat{\subsubsection}
  {\color{kudark}\normalsize\bfseries}{\thesubsubsection}{0.5em}{}
\titlespacing*{\section}{0pt}{1.4\baselineskip}{0.5\baselineskip}

% Minimum room a heading needs at the bottom of a page; if only ~1–3 lines are
% left (less than this), the heading is bumped to the next page instead of being
% stranded above its table/text.
\renewcommand{\bottomtitlespace}{4\baselineskip}

% ---- Hierarchical indentation by heading level ------------------------------
% Each heading level shifts BOTH the heading and the body text under it further
% to the right: section = 0, subsection (x.y) = 1 step, subsubsection (x.y.z) =
% 2 steps. We do this by having the sectioning commands set \leftskip (a global
% left margin that carries through every following line until the next heading
% resets it). \par first, so the *previous* paragraph is not re-indented.
% \leftskip indents plain paragraphs; \@totalleftmargin makes list environments
% (itemize/enumerate) start from the same indent instead of the page margin.
\makeatletter
\newlength{\lvlindent}
\setlength{\lvlindent}{1.8em}          % one indent step
\let\ulmsSecIndent\section
\renewcommand{\section}{\par\global\leftskip=0pt\global\@totalleftmargin=0pt\relax\ulmsSecIndent}
\let\ulmsSubsecIndent\subsection
\renewcommand{\subsection}{\par\global\leftskip=\lvlindent\global\@totalleftmargin=\lvlindent\relax\ulmsSubsecIndent}
\let\ulmsSubsubsecIndent\subsubsection
\renewcommand{\subsubsection}{\par\global\leftskip=2\lvlindent\global\@totalleftmargin=2\lvlindent\relax\ulmsSubsubsecIndent}
\makeatother

% Indent the headings themselves to match (titlesec ignores \leftskip for the
% heading line, so its left margin is set here via the first \titlespacing arg).
\titlespacing*{\subsection}   {\lvlindent} {1.0\baselineskip}{0.4\baselineskip}
\titlespacing*{\subsubsection}{2\lvlindent}{0.8\baselineskip}{0.3\baselineskip}

% ---- Lists (tighter, custom bullets) ----------------------------------------
\usepackage{enumitem}
\setlist[itemize]{leftmargin=1.6em, itemsep=2pt, topsep=3pt}
\setlist[enumerate]{leftmargin=1.9em, itemsep=2pt, topsep=3pt}

% ---- Tables ------------------------------------------------------------------
\usepackage{array}
\usepackage{booktabs}         % \toprule / \midrule / \bottomrule
\usepackage{tabularx}         % auto-width columns that wrap long Thai text
\usepackage{longtable}        % tables that can break across pages
\usepackage{multirow}         % merged cells for the team-structure table

% Left-aligned wrapping column types:
%   L{width} — fixed-width ragged-right paragraph column
%   Y        — flexible ragged-right X column (for tabularx)
\newcolumntype{L}[1]{>{\raggedright\arraybackslash}p{#1}}
\newcolumntype{Y}{>{\raggedright\arraybackslash}X}
\renewcommand{\arraystretch}{1.3}

% Helpers for a styled header row inside tables.
%   \hrow  — starts a header row and shades it (must be FIRST in the row;
%            colortbl allows only one \rowcolor per row).
%   \thead — bold header-cell text.
\newcommand{\hrow}{\rowcolor{headbg}}
\newcommand{\thead}[1]{\textbf{#1}}

% \begin{keeptable} … \end{keeptable} — keeps a table and the bold label above
% it on the SAME page (used by the data dictionary in §5). A tabularx cannot
% split across pages, so LaTeX is free to break between the label and its
% table, stranding the label alone at the bottom of a page; wrapping the pair
% in a minipage makes them one unbreakable block.
%   * \leftskip is zeroed inside: the minipage box is already placed at the
%     current heading indent, so leaving it set would indent the content twice
%     and push the table past the right margin.
%   * \parskip is restored because minipage's \@parboxrestore zeroes it, which
%     would glue the table straight onto the label line.
% Tables inside must use \linewidth (= the minipage width), not \textwidth.
\newenvironment{keeptable}
  {\par\medskip\noindent
   \begin{minipage}{\dimexpr\textwidth-\leftskip\relax}%
   \leftskip=0pt\relax
   \setlength{\parskip}{0.5\baselineskip plus 2pt}}
  {\end{minipage}\par\medskip}

% ---- Table of contents styling ----------------------------------------------
% By default LaTeX gives dotted leaders to subsections but NOT to sections
% (they print bold with just a right-aligned page number). tocloft lets us turn
% on dots for the section level too, so every TOC line runs "….." to its page.
\usepackage{tocloft}
\renewcommand{\cftsecdotsep}{\cftdotsep}   % dotted leaders for section entries
\renewcommand{\cftsecleader}{\cftdotfill{\cftsecdotsep}}
\renewcommand{\cftsecpagefont}{\normalfont}  % keep page numbers un-bold

% ---- Figures -----------------------------------------------------------------
\usepackage{graphicx}
\graphicspath{{figures/}}

% Landscape (rotated) full-page floats — needed by the ER diagram in §5, which
% is far wider than it is tall and would be unreadable scaled to \textwidth.
\usepackage{rotating}

% ---- Flowcharts / diagrams (TikZ) -------------------------------------------
% Used for the borrow–return workflow diagram (§5). Libraries: geometric shapes
% for decision diamonds, arrows.meta for arrowheads, positioning for relative
% node placement, calc/fit/backgrounds for routing loops and grouping.
\usepackage{tikz}
\usetikzlibrary{shapes.geometric, arrows.meta, positioning, calc, fit, backgrounds}

% ---- Links & bookmarks -------------------------------------------------------
\usepackage{hyperref}
\hypersetup{
  colorlinks=true, linkcolor=black, citecolor=linkblue, urlcolor=linkblue,
  pdflang={th},
}

% Thai labels for auto-generated headings.
\renewcommand{\contentsname}{สารบัญ (Table of Contents)}
\renewcommand{\tablename}{ตาราง}
\renewcommand{\figurename}{รูปที่}



\hypersetup{pdftitle={รายงานความก้าวหน้าโครงการ ULMs ครั้งที่ 5}}

% ---- ย่อหน้า (เหมือนฉบับที่ 4) ------------------------------------------------
\setlength{\parindent}{1.5em}
\setlength{\parskip}{0.35\baselineskip plus 2pt}
\titlespacing{\section} {0pt} {1.4\baselineskip}{0.5\baselineskip}
\titlespacing{\subsection} {\lvlindent} {1.0\baselineskip}{0.4\baselineskip}
\titlespacing{\subsubsection}{2\lvlindent}{0.8\baselineskip}{0.3\baselineskip}

\setlist[itemize]{leftmargin=1.6em, itemsep=2pt, topsep=0.55\baselineskip}
\setlist[enumerate]{leftmargin=1.9em, itemsep=2pt, topsep=0.55\baselineskip}

\setlength{\emergencystretch}{3em}
\newcommand{\zb}{\hspace{0pt}}

\usepackage{float}

% ---- ตัวแปรของรอบรายงาน ------------------------------------------------------
\newcommand{\reportdate}{24 กันยายน 2569}
\newcommand{\reportperiod}{18 – 24 กันยายน 2569}
\newcommand{\actualpct}{84\%}
\newcommand{\planpct}{94\%}
\newcommand{\prevactualpct}{79\%}

\begin{document}

% =============================================================================
% TITLE PAGE
% =============================================================================
\begin{titlepage}
\centering
\vspace*{2.8cm}

{\Large\bfseries\color{kugreen} รายงานความก้าวหน้าโครงการ ครั้งที่ 5\par}
\vspace{0.6cm}
{\LARGE\bfseries\color{kugreen} ระบบบริหารจัดการการยืม–คืนอุปกรณ์มหาวิทยาลัย\par}
\vspace{0.35cm}
{\Large\color{kudark} University Lending Management System (ULMs)\par}

\vspace{1.8cm}
\rule{0.72\textwidth}{0.8pt}
\vspace{0.9cm}

\begin{tabular}{@{}r@{\hspace{1.2em}}l@{}}
\textbf{รอบรายงาน} & \reportperiod \\[0.4em]
\textbf{วันที่รายงาน} & \reportdate \\[0.4em]
\textbf{ขอบเขตการนับ} & งานพัฒนาระบบ Sprint 0–5 (ไม่รวมงานก่อน 24 ก.ค.) \\[0.4em]
\textbf{สถานะ Sprint} & วันที่ 14 จาก 21 ของ Sprint 4 (Sprint ที่ 5 จาก 6) \\[0.4em]
\textbf{ความก้าวหน้า} & \actualpct\ เทียบแผน \planpct\ \ (ต่ำกว่าแผน 10\%) \\[0.4em]
\textbf{งบประมาณ} & ใช้ไป 0 บาท จากวงเงินประเมิน 5{,}000 บาท \\[0.4em]
\textbf{กำหนดเสร็จ} & 8 ตุลาคม 2569 (เลื่อนจาก 1 ตุลาคม 1 สัปดาห์) \\
\end{tabular}

\vspace{0.9cm}
\rule{0.72\textwidth}{0.8pt}

\vfill
{\large\color{kudark} คณะทำงาน ``miro ปั่น 14 บาท''\par}
\vspace{0.3cm}
{\color{kudark} ภาควิชาวิศวกรรมคอมพิวเตอร์ คณะวิศวกรรมศาสตร์\par}
{\color{kudark} มหาวิทยาลัยเกษตรศาสตร์\par}
\vspace{1.2cm}
\end{titlepage}

% =============================================================================
\section{บทสรุป}

รอบรายงานนี้มีการเปลี่ยนแปลงระดับแผนสามเรื่อง ซึ่งต้องอ่านก่อนตัวเลขทุกตัว

เรื่องแรก กำหนดเสร็จเลื่อนออกไปหนึ่งสัปดาห์ เป็น 8 ตุลาคม 2569 และผู้รับผิดชอบโครงการให้ยืด Sprint 4 ออกไปรับสัปดาห์นั้น วันรายงานจึงไม่ใช่วันปิด Sprint อีกต่อไป แต่เป็นวันที่ 14 จาก 21 ของ Sprint 4

เรื่องที่สอง สัปดาห์ที่เพิ่มมาเป็นเวลาไล่ตามงานที่ยังไม่เสร็จ ไม่ใช่การเลื่อนเป้าหมาย ค่าเป้าหมายรายวันในภาคผนวกจึงคงเดิมทุกช่อง แผน ณ วันรายงานยังเป็น \planpct\ และแผนยังแตะ 100\% ที่ 1 ตุลาคมเหมือนเดิม คณะทำงานเลือกวิธีนี้เพื่อไม่ให้ระยะห่างดูแคบลงเองจากการเลื่อนกำหนด ทั้งที่ปริมาณงานที่ทำได้ไม่ได้เปลี่ยน

เรื่องที่สาม ตะกร้าข้ามหน่วยงานและการจองล่วงหน้า ซึ่งรายงานครั้งที่ 4 ตัดออกตามแผนสำรอง ถูกดึงกลับเข้าขอบเขตที่ส่งมอบ เพราะมีเวลาเพิ่มขึ้นและฝั่งเซิร์ฟเวอร์กลับมาเดินแล้ว สองรายการนี้เหลืองานไม่เท่ากัน ตะกร้าข้ามหน่วยงานยังไม่มีโค้ดเลยทั้งสองฝั่ง ส่วนการจองล่วงหน้าหน้าเว็บทำไว้แล้ว โดยจำกัดการจองไม่เกิน 90 วัน แต่เซิร์ฟเวอร์ยังไม่บังคับกฎนั้น งานที่ต้องทำจึงโตขึ้น และค่าความก้าวหน้าของสองมิติฝั่งการพัฒนาลดลงตามตัวหารที่โตขึ้น

ความก้าวหน้าเพิ่มจาก \prevactualpct\ เป็น \actualpct\ ขณะที่แผนขยับจาก 86\% เป็น \planpct\ ระยะห่างจากแผนจึงกว้างขึ้นจาก 7\% เป็น 10\% ซึ่งเป็นระยะห่างที่กว้างที่สุดตั้งแต่เริ่มโครงการ สาเหตุไม่ใช่ว่าทีมทำงานได้น้อยลง รอบนี้เป็นรอบที่งานเข้ามากที่สุด แต่เป็นรอบแรกที่หักคุณภาพของหน้าจอออกจากค่าความก้าวหน้า และเป็นรอบที่ขอบเขตถูกขยายกลับ ทั้งสองอย่างดันตัวเลขลงพร้อมกัน

ฉบับนี้เปลี่ยนวิธีคิดความก้าวหน้าไปหนึ่งข้อ และต้องแจ้งไว้ก่อนอ่านตัวเลขอื่น รายงานครั้งที่ 1 ถึง 4 คิดจากการนับของที่สร้างเสร็จอย่างเดียว คือจำนวน procedure จำนวนหน้าจอ และจำนวนชุดทดสอบ ฉบับนี้หักข้อบกพร่องที่พบจากการตรวจหน้าจอด้วยมือเมื่อ 19 กันยายน ออกจากค่าความก้าวหน้าด้วย ถ้านับแบบเดิมล้วนจะได้ 90\% แต่เมื่อหักข้อบกพร่องที่ยังไม่แก้แล้วเหลือ \actualpct\ การตรวจนั้นเกิดขึ้นหลังรายงานครั้งที่ 4 จึงไม่เคยถูกหักมาก่อน คณะทำงานไม่ย้อนแก้ค่าของฉบับที่ส่งไปแล้ว รายละเอียดของการหักอยู่ในหัวข้อ \ref{sec:uireview}

รอบนี้เป็นรอบที่มีงานรวมเข้า main มากที่สุดตั้งแต่เริ่มโครงการ คือ 71 commit ผ่าน 23 pull request และต่างจากรอบก่อนตรงที่เป็นความสามารถใหม่จริง ไม่ใช่การแก้ความถูกต้อง จำนวน procedure ฝั่งเซิร์ฟเวอร์เพิ่มจาก 86 เป็น 102 และ domain เพิ่มจาก 10 เป็น 11

เรื่องสำคัญที่สุดคือการอุทธรณ์มีโค้ดฝั่งเซิร์ฟเวอร์แล้ว ซึ่งเป็นงานที่รายงานสามฉบับก่อนหน้าระบุตรงกันว่ายังไม่มีโค้ดเลย ตอนนี้เป็น domain ที่มี 6 procedure และหน้าอุทธรณ์ของอาจารย์เลิกใช้ข้อมูลจำลองแล้ว นอกจากนี้การจองห้องฝั่งเซิร์ฟเวอร์ใช้งานได้ หน้าส่งมอบของเจ้าหน้าที่และหน้าส่งออกรายงานซึ่งเคยเป็นหน้าว่างทั้งคู่ ถูกต่อเข้ากับเซิร์ฟเวอร์แล้ว ทั้งระบบจึงไม่เหลือหน้าว่างและไม่เหลือหน้าที่แสดงข้อมูลจำลองอีก

ความเสี่ยงเรื่องสัญญาระหว่างสองฝั่งที่รายงานครั้งที่ 4 แจ้งไว้ ปิดลงแล้วในเชิงตัวเลข ไฟล์สัญญาที่หน้าเว็บใช้ประกาศ 102 procedure ตรงกับเซิร์ฟเวอร์ทุกชื่อ ต่างกัน 0 รายการ จากเดิมที่ต่างกัน 14 รายการ หน้ารับของซึ่งรอบก่อนเรียกคำสั่งที่เซิร์ฟเวอร์ไม่มี จึงใช้งานได้แล้ว แต่ไฟล์นั้นยังเขียนด้วยมือ ความตรงกันนี้จึงมาจากการไล่แก้ ไม่ได้มาจากเครื่องมือที่กันไม่ให้ต่างกันตั้งแต่ต้น

เรื่องที่ต้องรายงานตรงไปตรงมาคือชุดทดสอบผ่านเบราว์เซอร์ รอบนี้เพิ่มจาก 13 เป็น 25 กรณี แต่มี 5 กรณีที่ไม่ผ่าน ทั้งห้ากรณีไม่ใช่ข้อผิดพลาดของเซิร์ฟเวอร์ แต่เป็นชุดทดสอบที่ค้นหาองค์ประกอบบนหน้าจอด้วยชื่อเดิม ขณะที่หน้าจอถูกแก้โดยตั้งใจในรอบนี้ สองในห้ากรณีคือชุดทดสอบการส่งคำขอยืม ซึ่งเป็นชุดเดียวกับที่เพิ่มเข้ามาในรอบก่อนเพื่อให้ผ่านจุดวัดข้อ 2 สาเหตุที่ไม่มีใครรู้ว่าชุดนี้เสียคือระบบตรวจอัตโนมัติยังไม่เรียกชุดทดสอบผ่านเบราว์เซอร์ ซึ่งคือจุดวัดข้อ 3 ที่ไม่ผ่านอีกรอบ

มิติที่ต้องเฝ้าระวังยังเป็น Test เหมือนสามรอบก่อน แม้ตัวเลขดิบจะขยับมาก คือชุดทดสอบฝั่ง backend เพิ่มจาก 266 เป็น 421 กรณี และชุดทดสอบส่วนประกอบฝั่งหน้าเว็บเพิ่มจาก 33 เป็น 105 กรณี

% =============================================================================
\section{แผนภาพรวมทั้ง 6 Sprint}
\label{sec:overview}

\noindent\begin{tabularx}{\dimexpr\textwidth-\leftskip\relax}{@{} L{2.9cm} Y L{2.2cm} @{}}
\toprule
\hrow \thead{Sprint} & \thead{ธีมและงานหลัก} & \thead{สถานะ} \\
\midrule
\textbf{Sprint 0} \newline 24 – 30 ก.ค. & วางฐาน — ปิดขอบเขต · ติดตั้งโครงสร้างพื้นฐานทั้งสองฝั่ง · database schema แกนกลาง · ระบบออกแบบและ mock & ผ่าน \\
\textbf{Sprint 1} \newline 31 ก.ค. – 13 ส.ค. & ล็อก contract และเปิดทาง — ล็อก contract API · ระบบเข้าสู่ระบบและออกจากระบบ · CI · เกณฑ์การยอมรับของวงจรหลัก & ผ่าน \\
\textbf{Sprint 2} \newline 14 – 27 ส.ค. & วงจรยืม–คืนของอุปกรณ์เคลื่อนย้ายได้ — รายการอุปกรณ์และจำนวนคงเหลือ · สร้างคำขอ อนุมัติ ส่งมอบ รับคืน · ข้อมูล seed · เลิกใช้ mock ในเส้นทางหลัก & ไม่ผ่านเกณฑ์ \newline (ยกงานไป Sprint 3) \\
\textbf{Sprint 3} \newline 28 ส.ค. – 10 ก.ย. & กฎธุรกิจรายระดับและการจองสถานที่ (T3) — อนุมัติตามระดับอุปกรณ์ · ผูกหมายเลขชิ้น · สล็อตเวลาและการจองห้อง · ตรวจสภาพและระบบเครดิต & ไม่ผ่านเกณฑ์ \newline (ยกงาน 4 วันไป Sprint 4) \\
\textbf{Sprint 4} \newline 11 ก.ย. – 1 ต.ค. \newline (ยืดจากเดิม 1 สัปดาห์) & ข้ามหน่วยงานและงานขัดเงา — หน้าจอที่เหลือ · การอุทธรณ์และการจองห้อง · ต่ออายุ \newline สี่วันแรกรับงานค้างของ Sprint 3 ก่อน \newline สัปดาห์ที่ยืดออกรับตะกร้าข้ามหน่วยงานและการจองล่วงหน้ากลับเข้ามา พร้อมข้อบกพร่องของหน้าจอที่ยังค้าง & กำลังทำ \newline (ผ่านไป 14 จาก 21 วัน) \\
\textbf{Sprint 5} \newline 2 – 8 ต.ค. & หยุดเพิ่มความสามารถ — ทดสอบถดถอย · นำขึ้นใช้งานจริง · คู่มือผู้ใช้และคู่มือติดตั้ง · ซ้อมนำเสนอ & ตามแผน \\
\bottomrule
\end{tabularx}

% =============================================================================
\section{สถานะโครงการโดยรวม}

\noindent\begin{tabularx}{\dimexpr\textwidth-\leftskip\relax}{@{} L{1.5cm} L{3.0cm} L{2.4cm} Y @{}}
\toprule
\hrow \thead{ด้าน} & \thead{แผน} & \thead{ผลจริง} & \thead{การประเมิน} \\
\midrule
เวลา & ผ่านไป 8.9 จาก 11 สัปดาห์ (81\%) & ผ่านไป 8.9 จาก 11 สัปดาห์ (81\%) & วันเสร็จเลื่อนจาก 1 ตุลาคม เป็น 8 ตุลาคม 2569 เหลือเวลาอีก 2 สัปดาห์ สัปดาห์แรกคือส่วนที่ยืดออกของ Sprint 4 ใช้ไล่ตามงานที่ค้างและงานที่ดึงกลับเข้ามา สัปดาห์ที่สองคือ Sprint 5 ทั้งช่วง สงวนไว้สำหรับการทดสอบและนำขึ้นใช้งาน ไม่ใช่การเพิ่มความสามารถ \\
งาน & \planpct & \actualpct & ต่ำกว่าแผน 10\% กว้างที่สุดตั้งแต่เริ่มโครงการ เป้าหมายไม่ได้เลื่อนตามกำหนดเสร็จ ค่าจริงถูกหักด้วยข้อบกพร่องของหน้าจอที่ยังไม่แก้ 9 ข้อ และถูกลดอีกจากการดึงงาน 2 รายการกลับเข้าขอบเขต รายละเอียดอยู่ในหัวข้อ \ref{sec:dimension} และ \ref{sec:uireview} \\
เงิน & 5{,}000 บาท & 0 บาท (0\%) & ยังไม่มีรายจ่าย รายการที่ประเมินไว้คือค่าเช่า Cloud/Hosting และฐานข้อมูลกลาง ซึ่งจะเกิดขึ้นเมื่อนำระบบขึ้นใช้งานจริงใน Sprint 5 \\
\bottomrule
\end{tabularx}

% =============================================================================
\clearpage
\section{ความก้าวหน้าแยกตามมิติของงาน}
\label{sec:dimension}

ค่าในคอลัมน์ ``จริง'' นับเฉพาะงานที่รวมเข้า main แล้ว ค่าในคอลัมน์ ``แผน'' อ่านจากแถว 24 กันยายนของภาคผนวกโดยตรง ซึ่งเป็นค่าเดิมที่ตั้งไว้ตั้งแต่ต้น ไม่ได้เลื่อนตามกำหนดเสร็จที่ขยับออกไป

คอลัมน์ ``นับล้วน'' คือค่าที่ได้จากการนับของที่สร้างเสร็จแบบเดียวกับรายงานสี่ฉบับก่อน โดยคิดบนขอบเขตที่ขยายกลับแล้ว ส่วนคอลัมน์ ``จริง'' คือค่าหลังหักข้อบกพร่องของหน้าจอที่ยังไม่แก้ตามหัวข้อ \ref{sec:uireview} สองมิติที่ถูกหักมากที่สุดคือ Front-End Dev และ Front-End Design เพราะข้อบกพร่องเกือบทั้งหมดอยู่ฝั่งหน้าจอ

\noindent\begin{tabularx}{\dimexpr\textwidth-\leftskip\relax}{@{} L{3.0cm} L{1.0cm} L{0.9cm} L{1.1cm} L{0.9cm} L{1.0cm} Y @{}}
\toprule
\hrow \thead{มิติงาน} & \thead{น้ำหนัก} & \thead{แผน} & \thead{นับล้วน} & \thead{จริง} & \thead{ผลต่าง} & \thead{ข้อเท็จจริงประกอบ} \\
\midrule
Back-End Design \newline (API + Database Model) & 10\% & 100\% & 100\% & 100\% & $0$ & contract เพิ่มเป็น 11 domain 102 procedure โดย domain ที่เพิ่มคือ appeal ซึ่งเป็นชิ้นสุดท้ายที่ยังไม่มีแบบ ตารางเพิ่มเป็น 33 และ migration เป็น 15 การออกแบบครบทุกเรื่องที่อยู่ในขอบเขตแล้ว \\
Front-End Design \newline (UI/UX + Business Process) & 10\% & 100\% & 100\% & 92\% & $-8$ & หน้าจอ 34 หน้ามีเนื้อหาจริงทั้งหมด ไม่เหลือหน้าว่าง แต่ผลตรวจหน้าจอยังค้างเรื่องระดับการออกแบบ 4 ข้อ คือไม่มีปฏิทินบอกวันที่ของเต็ม รายการอุปกรณ์ยังไม่ผูกกับช่วงเวลาที่ขอ ตารางบทลงโทษรวมการหักเครดิตกับการระงับสิทธิ์ไว้ด้วยกัน และใบจองห้องไม่มีรหัสครุภัณฑ์ \\
Back-End Dev & 30\% & 97\% & 90\% & 88\% & $-9$ & procedure เพิ่มจาก 86 เป็น 102 การอุทธรณ์มีโค้ดแล้ว 6 procedure การจองห้องใช้งานได้ การสมัครสมาชิกพร้อมยืนยันอีเมลใช้งานได้ cron job ทำงานจริงครบทั้ง 6 งาน ค่านับล้วนลดจาก 93 เป็น 90 เพราะขอบเขตที่ขยายกลับขาดฝั่งเซิร์ฟเวอร์ทั้งสองรายการ คือตะกร้าข้ามหน่วยงาน และการบังคับเพดานจองล่วงหน้า 90 วัน ยังเหลือการส่งอีเมลจริงด้วย \\
Front-End Dev & 25\% & 97\% & 90\% & 82\% & $-15$ & หน้าส่งมอบของและหน้าส่งออกรายงานต่อเข้าเซิร์ฟเวอร์แล้ว ทั้งคู่เคยเป็นหน้าว่าง หน้าห้องทั้ง 3 หน้าและหน้าอุทธรณ์อ่านจากเซิร์ฟเวอร์จริง ไม่เหลือหน้าที่แสดงข้อมูลจำลอง ค่านับล้วนลดจาก 92 เป็น 90 จากตะกร้าข้ามหน่วยงานที่ดึงกลับเข้าขอบเขต ส่วนการจองล่วงหน้าหน้าเว็บทำไว้แล้ว และถูกหักอีก 8 จุดจากข้อบกพร่องของหน้าจอที่ยังไม่แก้ 6 ข้อ กับที่แก้ได้บางส่วนอีก 4 ข้อ \\
Test (Front-End + Back-End) & 15\% & 80\% & 66\% & 62\% & $-18$ & ชุดทดสอบฝั่ง backend เพิ่มจาก 266 เป็น 421 กรณี ชุดส่วนประกอบฝั่งหน้าเว็บเพิ่มจาก 33 เป็น 105 กรณี service ที่มีชุดทดสอบเป็น 19 จาก 27 แต่ชุดผ่านเบราว์เซอร์ไม่ผ่าน 5 จาก 25 กรณี และระบบตรวจอัตโนมัติยังไม่เรียกชุดนั้น \\
Documentation & 10\% & 88\% & 88\% & 88\% & $0$ & SDS ปรับเป็นฉบับ 1.3 นับใหม่จากโค้ดของวันรายงาน สถานะที่เปลี่ยนมากที่สุดคือการอุทธรณ์ การจองห้อง และความตรงกันของสัญญาระหว่างสองฝั่ง \\
\midrule
\hrow \thead{รวมถ่วงน้ำหนัก} & \thead{100\%} & \thead{\planpct} & \thead{88\%} & \thead{\actualpct} & \thead{$-10$} & \\
\bottomrule
\end{tabularx}

มิติที่ต้องเฝ้าระวังยังเป็น Test ซึ่งห่างจากแผน 18 จุด ตัวเลขดิบขยับมากที่สุดในบรรดาทุกมิติ คือชุดทดสอบฝั่ง backend เพิ่ม 155 กรณีและชุดส่วนประกอบเพิ่ม 72 กรณี แต่สิ่งที่ยังขาดเป็นเรื่องเดิมสองข้อ คือชุดทดสอบของวงจรยืม–คืนที่เดินครบเส้นทาง ซึ่งปัจจุบันยังครอบคลุมเฉพาะขั้นส่งคำขอ และการเรียกชุดทดสอบผ่านเบราว์เซอร์จากระบบตรวจอัตโนมัติ ซึ่งไฟล์ตั้งค่ามีอยู่ในรีโปมาสามรอบแล้วแต่ยังไม่มีขั้นตอนใดเรียกใช้ เหตุผลที่ Test ถูกหักเพิ่มอีก 4 จุดคือข้อบกพร่องทั้ง 30 ข้อในหัวข้อ \ref{sec:uireview} ถูกพบด้วยการเปิดหน้าจอดูทีละหน้า ไม่มีข้อใดถูกจับโดยชุดทดสอบที่มีอยู่

มิติที่ทำงานได้มากที่สุดคือ Back-End Dev ซึ่ง procedure เพิ่ม 16 รายการ และงานที่เพิ่มคือสองเรื่องที่ค้างมานานที่สุด คือการอุทธรณ์และการจองห้อง แต่ค่าที่รายงานกลับห่างจากแผน 9 จุด เท่ากับรอบก่อน เพราะสองเหตุผลที่ไม่เกี่ยวกับปริมาณงาน เหตุผลแรกคือขอบเขตถูกขยายกลับ ตะกร้าข้ามหน่วยงานและการจองล่วงหน้ากลับมาเป็นงานที่ต้องส่งมอบ และทั้งสองยังไม่มีโค้ด เหตุผลที่สองคือข้อมูลที่หน้าจอต้องใช้แต่เซิร์ฟเวอร์ยังไม่ให้ คือเวลาที่ของจะว่างครั้งถัดไป และรหัสครุภัณฑ์บนใบจองห้อง

ข้อสังเกตที่คณะทำงานอยากให้อ่านคู่กับตารางนี้คือ ตัวเลขที่ลดลงสามตัวในรอบนี้ ไม่มีตัวไหนเกิดจากการทำงานได้น้อยลง ตัวแรกมาจากการหักคุณภาพของหน้าจอ ตัวที่สองมาจากขอบเขตที่ขยายกลับ ตัวที่สามมาจากการไม่เลื่อนเป้าหมายตามกำหนดเสร็จ ทั้งสามเป็นการเปลี่ยนวิธีวัดหรือเปลี่ยนสิ่งที่ต้องส่งมอบ ไม่ใช่การเปลี่ยนผลงาน

\begin{figure}[H]
\centering
\begin{tikzpicture}[x=0.062cm, y=1cm]
 \foreach \x in {0,20,40,60,80,100} {
 \draw[gray!25] (\x,-0.30) -- (\x,6.15);
 \node[below, font=\scriptsize] at (\x,-0.30) {\x};
 }
 \node[below, font=\scriptsize] at (50,-0.68) {\% ความก้าวหน้าของมิตินั้น};
 \draw[thick] (0,-0.30) -- (0,6.15);

 \foreach \y/\name/\wt/\plan/\act in {%
 5.6/{Back-End Design}/10/100/100,
 4.6/{Front-End Design}/10/100/92,
 3.6/{Back-End Dev}/30/97/88,
 2.6/{Front-End Dev}/25/97/82,
 1.6/{Test}/15/80/62,
 0.6/{Documentation}/10/88/88} {
 \node[left, xshift=-0.16cm, font=\scriptsize, align=right] at (0,\y)
 {\name\\[-3pt]{\tiny น้ำหนัก \wt\%}};
 \fill[gray!45] (0,{\y+0.05}) rectangle (\plan,{\y+0.30});
 \fill[kugreen] (0,{\y-0.30}) rectangle (\act,{\y-0.05});
 \node[right, font=\tiny, gray!45!black] at (\plan,{\y+0.175}) {\plan};
 \node[right, font=\tiny, kugreen] at (\act,{\y-0.175}) {\act};
 }

 \fill[gray!45] (0,-1.36) rectangle (7,-1.54);
 \node[right, font=\scriptsize] at (9,-1.45) {แผน};
 \fill[kugreen] (0,-1.76) rectangle (7,-1.94);
 \node[right, font=\scriptsize] at (9,-1.85) {ผลจริง};
\end{tikzpicture}
\vspace{0.6em}
\caption{ความก้าวหน้ารายมิติ ณ \reportdate\ ซึ่งเป็นวันที่ 14 จาก 21 ของ Sprint 4
 เทียบกับค่าเป้าหมายเดิมที่แถว 24 กันยายนของภาคผนวก ซึ่งไม่ได้เลื่อนตามกำหนดเสร็จ
 ตัวเลขใต้ชื่อมิติคือน้ำหนักที่ใช้ถ่วงเมื่อรวมเป็นค่าเดียว
 แท่งผลจริงนับเฉพาะงานที่รวมเข้า main แล้ว
 แท่งผลจริงเป็นค่าหลังหักข้อบกพร่องของหน้าจอที่ยังไม่แก้แล้ว
 มิติที่ห่างจากแผนมากที่สุดคือ Test ที่ $-18$ รองลงมาคือ Front-End Dev ที่ $-15$
 ส่วน Back-End Design และ Documentation เดินตรงตามแผนพอดี}
\label{fig:dims}
\end{figure}

% =============================================================================
\clearpage
\section{กราฟความก้าวหน้าเทียบแผน}

เส้นแผนสร้างจากแผน Sprint โดยตรง คือกำหนดค่าเป้าหมายรายมิติ ณ วันสิ้นสุดของแต่ละ Sprint แล้วคูณน้ำหนักรวมกันเป็นค่าสะสม ตารางค่าทั้งหมดอยู่ในภาคผนวก แกนนอนเริ่มที่ 24 กรกฎาคม 2569 ซึ่งเป็นวันเริ่มลงมือพัฒนา

ความชันของเส้นผลจริงในรอบนี้คือ 5 จุดในเจ็ดวัน ขณะที่เส้นแผนชัน 16 จุดในช่วง Sprint 4 ระยะห่างจึงกว้างขึ้นจาก 7 จุดเป็น 10 จุด ซึ่งกว้างที่สุดตั้งแต่เริ่มโครงการ

เส้นแผนไม่ถูกเลื่อนตามกำหนดเสร็จที่ขยับออกไป ค่าเป้าหมายทุกจุดยังอยู่วันเดิม และยังแตะ 100\% ที่ 1 ตุลาคมเหมือนที่ประกาศไว้ตั้งแต่ต้น ช่วง 1 ถึง 8 ตุลาคมเส้นแผนจึงราบอยู่ที่ 100 เพราะไม่มีเป้าหมายใหม่ในช่วงนั้น มีแต่เวลาให้เส้นผลจริงไล่ขึ้นไปหา ระยะในแนวตั้งระหว่างสองเส้น ณ 1 ตุลาคม คือปริมาณงานที่ต้องปิดให้ได้ภายในสัปดาห์สุดท้าย

ถ้านับแบบสี่ฉบับก่อนโดยไม่หักข้อบกพร่องของหน้าจอ จุดของวันนี้จะอยู่ที่ 88\% วงกลมกลวงในรูปคือค่านั้น

\begin{figure}[H]
\centering
\begin{tikzpicture}[x=1.20cm, y=0.058cm]
 \foreach \y in {0,20,40,60,80,100}
 \draw[gray!22] (0,\y) -- (11.1,\y);

 \foreach \x in {0.86, 2.86, 4.86, 6.86, 9.86}
 \draw[gray!40, dashed, thin] (\x,0) -- (\x,108);

 \foreach \x/\lbl in {0.43/{S0}, 1.86/{S1}, 3.86/{S2}, 5.86/{S3}, 8.36/{S4}, 10.36/{S5}}
 \node[font=\scriptsize\bfseries, gray!50!black] at (\x,114) {\lbl};

 \draw[->, thick] (0,0) -- (11.5,0);
 \draw[->, thick] (0,0) -- (0,122) node[above, font=\scriptsize] {\% ความก้าวหน้าสะสม};
 \foreach \y in {0,20,40,60,80,100}
 \draw (0,\y) -- (-0.10,\y) node[left, font=\scriptsize] {\y};

 \foreach \x/\lbl in {0/{24 ก.ค.}, 0.86/{30 ก.ค.}, 2.86/{13 ส.ค.},
 4.86/{27 ส.ค.}, 6.86/{10 ก.ย.}, 8.86/{24 ก.ย.}, 9.86/{1 ต.ค.}, 10.86/{8 ต.ค.}}
 {\draw (\x,0) -- (\x,-3);
 \node[rotate=45, anchor=north east, font=\scriptsize] at (\x,-4) {\lbl};}

 % ---- เส้นแผน ----
 \draw[thick, kudark, densely dashed]
 plot coordinates {(0,1) (0.86,12) (2.86,42) (4.86,65) (6.86,78) (8.86,94) (9.86,100) (10.86,100)};
 \foreach \p in {(0,1), (0.86,12), (2.86,42), (4.86,65), (6.86,78), (8.86,94), (9.86,100)}
 \fill[kudark] \p circle (1.6pt);

 % ---- เส้นผลจริง — จุดสุดท้าย 8.86 = 24 ก.ย. = วันรายงานนี้ ----
 \draw[very thick, kugreen]
 plot coordinates {(0,0) (2.0,24) (2.43,32) (2.86,43) (4.29,43) (4.86,52) (6.0,70) (6.86,75) (7.86,79) (8.86,84)};
 \foreach \p in {(0,0), (2.0,24), (2.43,32), (2.86,43), (4.29,43), (4.86,52), (6.0,70), (6.86,75), (7.86,79), (8.86,84)}
 \fill[kugreen] \p circle (2.2pt);

 % เส้นบางคือค่าที่ได้ถ้านับแบบเดิม ไม่หักข้อบกพร่องของหน้าจอ
 \draw[thin, kugreen, dotted] (7.86,79) -- (8.86,88);
 \fill[white] (8.86,88) circle (2.2pt);
 \draw[kugreen, thin] (8.86,88) circle (2.2pt);
 % ป้ายสองอันนี้ต้องอยู่ขวาของ x=8.86 เท่านั้น ที่ว่างจริงมีแค่ตรงนั้น
 % ถ้าวางทางซ้ายแล้วถมขาว จะลบเส้นแผนที่ไต่ผ่านช่วง 85 ถึง 93 พอดี
 \node[anchor=west, font=\tiny, kugreen] at (8.95,86) {นับล้วน 88};
 \draw[gray!70, thick, |-|] (8.86,84) -- (8.86,94);
 \node[anchor=west, font=\tiny, gray!70!black] at (8.95,91) {$-10$};

 % ไม่มีป้ายกำกับจุดของรอบก่อน (17 ก.ย.) บนกราฟนี้โดยตั้งใจ
 % ช่องว่างระหว่างเส้นแผนกับเส้นผลจริงตรงจุดนั้นกว้างแค่ 2.6 ถึง 7 หน่วย
 % แต่ป้ายสองบรรทัดต้องการราว 10 หน่วย วางชิดจุดเมื่อไรก็ทับเส้น
 % ทางเดียวที่เหลือคือลากเส้นชี้ ซึ่งทำให้จุดของรอบก่อนเด่นแข่งกับวันรายงาน
 % ค่าของรอบก่อนอยู่ในหัวข้อ 1 และในภาคผนวกอยู่แล้ว จึงไม่ต้องซ้ำบนกราฟ
 % เส้นชี้จากจุดวันรายงานลงมาที่ป้าย เพื่อไม่ให้ป้ายลอยอยู่กลางที่ว่าง
 \draw[gray!55, thin] (8.86,82) -- (8.60,72);
 \node[anchor=north west, font=\scriptsize, kugreen, align=left,
 fill=white, inner sep=1.5pt] at (8.20,71)
 {24 ก.ย. (วันรายงาน)\\ผลจริง 84\% · แผน 94\%};

 \fill[white] (0.55,72) rectangle (4.60,96);
 \draw[thick, kudark, densely dashed] (0.70,89) -- (1.40,89);
 \node[right, font=\scriptsize] at (1.50,89) {แผน (Planned)};
 \draw[very thick, kugreen] (0.70,79) -- (1.40,79);
 \node[right, font=\scriptsize] at (1.50,79) {ผลจริง (Actual)};
\end{tikzpicture}
\caption{กราฟ เปรียบเทียบความก้าวหน้าตามแผนกับผลการดำเนินงานจริง
 แกนนอนวางตามขอบ Sprint ทั้ง 6 ช่วง (เส้นประแนวตั้ง)
 จุดบนเส้นผลจริงคือวันที่วัดความก้าวหน้า
 วันรายงาน 24 กันยายน 2569 ตรงกับขอบ Sprint 4 พอดี
 ค่าแผน ณ วันนั้นจึงอ่านจากแถว S4 ของภาคผนวกโดยตรง
 เส้นประแนวตั้งคือขอบ Sprint หลังยืด Sprint 4 ออกไปถึง 1 ตุลาคม
 วงกลมกลวงที่ 88 คือค่าที่ได้ถ้านับแบบสี่ฉบับก่อน คือไม่หักข้อบกพร่องของหน้าจอ
 จุดทึบที่ 84 คือค่าที่ฉบับนี้รายงาน
 เส้นแผนราบที่ 100 ในช่วง 1 ถึง 8 ตุลาคม เพราะเป้าหมายไม่ได้เลื่อนตามกำหนดเสร็จ
 ช่วงนั้นเป็นเวลาที่ให้ไว้ไล่ตามงานที่ยังไม่เสร็จ
 จุดก่อนหน้า 24 กันยายน ทั้งหมดยังเป็นค่าที่ส่งไปแล้ว ไม่ได้คิดใหม่ด้วยวิธีของฉบับนี้
 ช่วงราบระหว่าง 13 ถึง 23 สิงหาคม คือช่วงสอบกลางภาค}
\label{fig:scurve}
\end{figure}

% =============================================================================
\clearpage
\section{ผลของจุดวัดที่รายงานครั้งที่ 4 ตั้งไว้}
\label{sec:review}

รายงานครั้งที่ 4 ประกาศจุดวัดไว้สี่ข้อ โดยแยกเป็นข้อละหนึ่งงานตามข้อสังเกตของรอบนั้น ตารางนี้คือผลของทั้งสี่ข้อ

\noindent\begin{tabularx}{\dimexpr\textwidth-\leftskip\relax}{@{} Y L{1.7cm} L{5.4cm} L{1.9cm} @{}}
\toprule
\hrow \thead{จุดวัด} & \thead{กำหนด} & \thead{สภาพจริง ณ วันรายงาน} & \thead{ผล} \\
\midrule
หน้าส่งมอบของเจ้าหน้าที่ใช้งานได้จริง & 19 ก.ย. & หน้าจอเรียก loan.getForStaff และ loan.swapUnit ได้จริง ไม่ใช่หน้าว่างอีกต่อไป แต่รวมเข้า main วันที่ 21 ก.ย. ช้ากว่ากำหนด 2 วัน & ผ่าน \newline (ช้ากว่ากำหนด) \\
วงจรยืม–คืนเดินครบเส้นทางปกติผ่านหน้าจอ ตั้งแต่ส่งคำขอถึงตรวจสภาพ & 21 ก.ย. & ชุดทดสอบผ่านเบราว์เซอร์ยังครอบคลุมเฉพาะขั้นส่งคำขอเหมือนรอบก่อน ไม่มีกรณีของขั้นอนุมัติ ส่งมอบ รับคืน หรือตรวจสภาพ และกรณีของการส่งคำขอเองก็ไม่ผ่านแล้ว 2 จาก 3 & ไม่ผ่าน \\
ระบบตรวจอัตโนมัติเรียกชุดทดสอบผ่านเบราว์เซอร์ และขวางการรวมโค้ดเมื่อไม่ผ่าน & 21 ก.ย. & ไฟล์ .github/workflows/ci.yml ไม่ถูกแก้เลยตลอดรอบ ยังไม่มีขั้นตอนใดเรียก playwright และไม่มีขั้นตอน lint ด้วย & ไม่ผ่าน \\
การจองห้องฝั่ง backend ใช้งานได้ หน้าห้องทั้ง 3 หน้าเลิกใช้ข้อมูลจำลอง และไม่เหลือหน้าว่างในระบบ & 24 ก.ย. & loan.createRoomBooking กับ item.roomAvailability ใช้งานได้ หน้าห้องทั้ง 3 หน้าอ่านจากเซิร์ฟเวอร์จริง และไม่มีไฟล์ใดใช้ PlaceholderPage แล้ว & ผ่าน \\
\bottomrule
\end{tabularx}

ผลรอบนี้คือผ่าน 2 ไม่ผ่าน 2 ดีขึ้นจากสองรอบก่อนที่ไม่ผ่าน 3 ข้อเท่ากัน ข้อที่ผ่านคือข้อที่เป็นงานพัฒนาความสามารถ ส่วนข้อที่ไม่ผ่านทั้งสองข้อเป็นเรื่องของการทดสอบ ซึ่งเป็นมิติเดียวกับที่ห่างจากแผนมากที่สุดมาตลอดสี่รอบ

ข้อสังเกตที่คณะทำงานบันทึกไว้คือจุดวัดข้อ 2 และข้อ 3 เกี่ยวพันกัน การที่ระบบตรวจอัตโนมัติไม่เรียกชุดทดสอบผ่านเบราว์เซอร์ ทำให้ไม่มีใครรู้ว่าชุดทดสอบของการส่งคำขอยืมที่เพิ่งเพิ่มในรอบก่อน เสียไปตั้งแต่กลางรอบนี้ รายละเอียดอยู่ในหัวข้อ \ref{sec:evidence}

% =============================================================================
\clearpage
\section{ผลตรวจหน้าจอด้วยมือ และการหักความก้าวหน้า}
\label{sec:uireview}

เมื่อวันที่ 19 กันยายน คณะทำงานเปิดระบบขึ้นมาตรวจทีละหน้าจนครบทั้งสี่บทบาท แล้วบันทึกทุกจุดที่ผิดหรือน่าสงสัยเป็นเอกสารภาพ 20 หน้า รวม 33 ข้อ แยกเป็นข้อบกพร่อง 30 ข้อ และคำถามเชิงออกแบบอีก 3 ข้อ ซึ่งไม่นับเป็นข้อบกพร่อง

การตรวจครั้งนั้นเกิดขึ้นหลังรายงานครั้งที่ 4 สองวัน ค่าความก้าวหน้าที่ส่งไปในฉบับนั้นจึงยังไม่ได้หักข้อบกพร่องเหล่านี้ ฉบับนี้เป็นฉบับแรกที่หัก

\subsection{สถานะของข้อบกพร่องทั้ง 30 ข้อ ณ วันรายงาน}

\noindent\begin{tabularx}{\dimexpr\textwidth-\leftskip\relax}{@{} L{4.8cm} L{1.6cm} Y @{}}
\toprule
\hrow \thead{สถานะ} & \thead{จำนวน} & \thead{ที่มาของการตัดสิน} \\
\midrule
แก้แล้ว ทีมกำกับเอง & 8 & มีป้าย Fixed หรือ Removed เขียนกำกับไว้ในเอกสารภาพ \\
แก้แล้ว ยังไม่มีป้ายกำกับ & 8 & คณะทำงานตรวจซ้ำจากโค้ดบน main ที่ e0a4865 แล้วยืนยันว่าแก้จริง \\
แก้บางส่วน & 5 & มีป้ายกำกับว่าแก้เบื้องต้น แก้เฉพาะเคอร์เซอร์ แก้เฉพาะฝั่งหน้าเว็บ หรือกำกับว่าไม่แน่ใจ \\
\hrow \thead{ยังไม่แก้} & \thead{9} & \thead{ไม่มีป้ายกำกับ และตรวจจากโค้ดแล้วยังไม่พบการแก้} \\
\bottomrule
\end{tabularx}

\subsection{ข้อที่แก้แล้วแต่ยังไม่มีใครกำกับ}

แปดข้อนี้ทีมแก้ไปแล้วระหว่างรอบ แต่ยังไม่ได้กลับไปกำกับในเอกสารภาพ คณะทำงานตรวจซ้ำจากโค้ดจึงนับเป็นแก้แล้ว

\begin{itemize}
\item หน้าส่งมอบของเจ้าหน้าที่ซึ่งเคยเป็นหน้าว่าง ต่อเข้าเซิร์ฟเวอร์แล้ว
\item หน้าส่งออกรายงานซึ่งเคยเป็นหน้าว่าง ดาวน์โหลดไฟล์ CSV ได้แล้ว
\item รายการห้อง 8 ห้องซึ่งเคยเป็นข้อมูลจำลอง อ่านจากเซิร์ฟเวอร์จริงแล้ว
\item หน้าอุทธรณ์ของอาจารย์เลิกใช้ข้อมูลตัวอย่าง และรูปก่อน-หลังที่เคยโหลดไม่ขึ้น แสดงได้แล้ว
\item จำนวนของคงเหลือที่ผู้ยืมกับเจ้าหน้าที่เคยเห็นไม่ตรงกัน ใช้กฎเดียวกันแล้ว
\item อุปกรณ์ระดับ T1 ไม่ต้องให้เจ้าหน้าที่พิมพ์หมายเลขเองอีก
\item ป้าย ``ถูกระงับการยืม'' ที่เคยขึ้นผิดกับคนที่ยังยืมได้ ถูกถอดออกแล้ว
\end{itemize}

\subsection{เก้าข้อที่ยังไม่แก้}

\noindent\begin{tabularx}{\dimexpr\textwidth-\leftskip\relax}{@{} L{3.4cm} Y L{2.4cm} @{}}
\toprule
\hrow \thead{หน้าจอ} & \thead{อาการ} & \thead{หักเข้ามิติ} \\
\midrule
รายการอุปกรณ์ของผู้ยืม & ตัวเลือกเรียง ``ยอดนิยม'' เรียงตามจำนวนชิ้นที่มี ไม่ใช่ความนิยม ตรวจในโค้ดแล้วยังเรียงด้วยผลต่างของจำนวนชิ้นทั้งหมด & Front-End Dev \\
ตะกร้าและการส่งคำขอ & ไม่มีตัวเลือกว่าต้องได้ครบทั้งตะกร้าหรือไม่เอาเลย ของที่ต้องใช้คู่กันจึงถูกส่งไปเฉพาะชิ้นที่ว่าง แล้วผู้ยืมต้องตามยกเลิกเอง & Front-End Design \\
รายการคำขอของผู้ยืม & แถบความคืบหน้าเดินเลยขั้นยืนยันการจองไปแล้ว แต่ป้ายสถานะยังเป็นรออนุมัติ & Front-End Dev \\
ใบจองห้อง & ไม่มีรหัสครุภัณฑ์ จึงเหลือขีดค้างไว้ในช่องนั้น & Back-End Dev \\
คลังอุปกรณ์ของเจ้าหน้าที่ & คอลัมน์ ``ว่างครั้งถัดไป'' เป็นขีดทุกแถว ค่านี้มีที่เก็บในสัญญาแล้วแต่เซิร์ฟเวอร์ยังไม่ส่งให้หน้าคลัง & Back-End Dev \\
คลังอุปกรณ์ของเจ้าหน้าที่ & ทั้งหน้าไม่มีปุ่มเพิ่มอุปกรณ์ มีแต่ช่องค้นหา เจ้าหน้าที่จึงลงทะเบียนของใหม่จากหน้านี้ไม่ได้ & Front-End Dev \\
หน้าส่งมอบของ & แถบนำทางเขียนว่าคิวงาน แล้วตามด้วยภาพรวมของฉัน ซึ่งไม่ตรงกับหัวข้อหน้าที่เป็นส่งมอบและรับคืน & Front-End Dev \\
ผู้ใช้ในหน่วยงาน & ตารางบทลงโทษรวมการหักเครดิตกับการระงับสิทธิ์ไว้ด้วยกัน ทั้งที่เป็นคนละเรื่องและมีผลต่างกัน & Front-End Design \\
รายการห้อง & แถบแจ้งเตือนอ้างถึงห้อง 7603 ซึ่งไม่ปรากฏในตารางด้านล่าง & Front-End Dev \\
\bottomrule
\end{tabularx}

\subsection{การหักคิดอย่างไร}

คณะทำงานหักเฉพาะมิติที่ข้อบกพร่องนั้นอยู่จริง ไม่หักรวมทั้งโครงการ ข้อบกพร่องเกือบทั้งหมดอยู่ฝั่งหน้าจอ สองมิติของฝั่งนั้นจึงรับการหักมากที่สุด

\noindent\begin{tabularx}{\dimexpr\textwidth-\leftskip\relax}{@{} L{3.2cm} L{1.5cm} L{1.4cm} L{1.2cm} Y @{}}
\toprule
\hrow \thead{มิติงาน} & \thead{นับล้วน} & \thead{หลังหัก} & \thead{หัก} & \thead{เหตุผล} \\
\midrule
Back-End Design & 100\% & 100\% & $0$ & สัญญาและโครงสร้างฐานข้อมูลไม่มีข้อบกพร่องในผลตรวจ \\
Front-End Design & 100\% & 92\% & $-8$ & 4 ข้อที่เป็นระดับการออกแบบ ไม่ใช่ความผิดพลาดของโค้ด \\
Back-End Dev & 90\% & 88\% & $-2$ & 2 ข้อที่หน้าจอต้องการข้อมูลแต่เซิร์ฟเวอร์ยังไม่ส่งให้ \\
Front-End Dev & 90\% & 82\% & $-8$ & 6 ข้อที่ยังไม่แก้ และ 4 ข้อที่แก้ได้บางส่วน \\
Test & 66\% & 62\% & $-4$ & ข้อบกพร่องทั้ง 30 ข้อถูกพบด้วยการเปิดดูทีละหน้า ไม่มีข้อใดถูกชุดทดสอบจับได้ \\
Documentation & 88\% & 88\% & $0$ & เอกสารบันทึกสถานะตามจริง รวมข้อที่ยังไม่แก้ \\
\midrule
\hrow \thead{รวมถ่วงน้ำหนัก} & \thead{88\%} & \thead{\actualpct} & \thead{$-4$} & \\
\bottomrule
\end{tabularx}

ข้อสังเกตที่คณะทำงานบันทึกไว้คือ สิบหกข้อจากสามสิบข้อแก้ไปแล้วภายในห้าวัน ซึ่งเร็ว แต่แปดข้อในนั้นไม่มีใครกลับไปกำกับสถานะในเอกสารภาพ ถ้าไม่ได้ตรวจซ้ำจากโค้ด รายงานฉบับนี้จะหักเกินจริงไปแปดข้อ การกำกับสถานะกลับไปที่เอกสารตรวจจึงควรเป็นส่วนหนึ่งของการปิดงานแต่ละข้อ ไม่ใช่งานที่ทำตอนท้าย

% =============================================================================
\clearpage
\section{ผลงานที่ส่งมอบในรอบนี้}

รอบนี้มี 71 commit รวมเข้า main ผ่าน 23 pull request จากผู้ร่วมงาน 8 คน เป็นรอบที่มีงานรวมเข้า main มากที่สุดตั้งแต่เริ่มโครงการ

\subsection{ฝั่ง backend}

\begin{itemize}
\item การอุทธรณ์มีโค้ดแล้วเป็น domain ที่ 11 มี 6 procedure ครอบคลุมการยื่น การอ่านรายการของตนเอง การอ่านรายการของผู้พิจารณา และการตัดสิน โดยผูกกับบทลงโทษจริงในระบบเครดิต
\item การจองห้องใช้งานได้ ทั้งการอ่านช่วงเวลาว่างและการจอง โดยบังคับกฎสล็อตเวลา เพดาน 3 ชั่วโมง และกฎวันเดียวกันที่ฝั่งเซิร์ฟเวอร์ทุกเส้นทาง
\item การสมัครสมาชิกด้วยตนเองพร้อมยืนยันอีเมล บัญชีที่สมัครใหม่ยังไม่เปิดใช้จนกว่าจะกดลิงก์ยืนยัน และคำตอบไม่บอกว่าอีเมลนั้นมีอยู่แล้วหรือไม่ พร้อมการขอและตั้งรหัสผ่านใหม่
\item เหตุผลการปฏิเสธเก็บในคอลัมน์ของตนเอง ไม่ทับเหตุผลที่ผู้ยืมกรอก โดยเพิ่มคอลัมน์และ migration รองรับ
\item การบันทึกร่องรอยการใช้งานเสร็จครบ รวมถึงการตัดสินคำขอต่ออายุ
\item cron job ทำงานจริงครบทั้ง 6 งาน โดยถอด computeAvailability และ rollupDailyStats ออกโดยตั้งใจ เพราะของคงเหลือคำนวณสดอยู่แล้ว
\end{itemize}

\subsection{ฝั่ง frontend}

\begin{itemize}
\item หน้าส่งมอบของเจ้าหน้าที่ต่อเข้าเซิร์ฟเวอร์แล้ว รวมถึงการเปลี่ยนชิ้นและการตรวจสิทธิ์ก่อนส่งมอบ
\item หน้าส่งออกรายงานดาวน์โหลดไฟล์ CSV จากข้อมูลจริง ทั้งสองหน้านี้เคยเป็นหน้าว่างในรายงานครั้งที่ 4
\item หน้าห้องทั้ง 3 หน้าอ่านจากเซิร์ฟเวอร์จริง และหน้าอุทธรณ์ของอาจารย์เลิกใช้ข้อมูลจำลอง
\item การขอต่ออายุส่งขึ้นเซิร์ฟเวอร์ ไม่เก็บไว้ในเบราว์เซอร์อีก
\item การคืนของต้องมีรูปหลังใช้งาน โดยเพิ่มขั้นถ่ายรูปในคิวงานของเจ้าหน้าที่
\item หน้าประวัติการซ่อมเป็นหน้าใหม่ พร้อมหน้าสมัครสมาชิก ยืนยันอีเมล และตั้งรหัสผ่านใหม่
\item ไฟล์สัญญาที่หน้าเว็บใช้ ประกาศ 102 procedure ตรงกับเซิร์ฟเวอร์ทุกชื่อ
\end{itemize}

\subsection{เครื่องมือและกระบวนการ}

\begin{itemize}
\item ชุดทดสอบฝั่ง backend เพิ่มจาก 266 เป็น 421 กรณี และจาก 24 เป็น 41 ชุด
\item ชุดทดสอบส่วนประกอบฝั่งหน้าเว็บเพิ่มจาก 33 เป็น 105 กรณี และจาก 3 เป็น 17 ไฟล์ ครอบคลุมการเข้าสู่ระบบ การยกเลิกคำขอ หน้าผู้ดูแลระบบ และคุณลักษณะที่ไม่ใช่หน้าที่ของระบบ
\item ชุดทดสอบผ่านเบราว์เซอร์เพิ่มจาก 13 เป็น 25 กรณี โดยเพิ่มชุดของหน้าผู้ดูแลระบบ
\item ระบบตรวจอัตโนมัติไม่ถูกแก้เลยในรอบนี้
\end{itemize}

\subsection{เอกสารและการออกแบบ}

\begin{itemize}
\item SDS ปรับเป็นฉบับ 1.3 นับสถานะใหม่จากโค้ดของวันรายงาน
\item สถานะที่เปลี่ยนมากที่สุดคือการอุทธรณ์และการจองห้อง ซึ่งเปลี่ยนจาก ``ยังไม่มีฟังก์ชัน'' เป็นใช้งานได้ และความต่างของสัญญาระหว่างสองฝั่งซึ่งเปลี่ยนจาก 14 รายการเป็น 0
\end{itemize}

% =============================================================================
\clearpage
\section{หลักฐานการทดสอบที่รันจริงในรอบนี้}
\label{sec:evidence}

ตัวเลขในตารางนี้ได้จากการรันบนเครื่องของคณะทำงานในวันรายงาน โดยดึงโค้ดจาก main ที่ e0a4865 สร้างฐานข้อมูลใหม่แล้ว seed ข้อมูลตั้งต้น จากนั้นสตาร์ต backend และ frontend เองก่อนรันชุดทดสอบผ่านเบราว์เซอร์ ไม่ได้อ่านจากผลของระบบตรวจอัตโนมัติ

\noindent\begin{tabularx}{\dimexpr\textwidth-\leftskip\relax}{@{} L{4.6cm} L{2.6cm} Y @{}}
\toprule
\hrow \thead{ชุดที่รัน} & \thead{ผล} & \thead{หมายเหตุ} \\
\midrule
ชุดทดสอบฝั่ง backend & ผ่าน 421 จาก 421 & 41 ชุดทดสอบ รายงานครั้งที่ 4 แจ้ง 266 กรณี ส่วนที่เพิ่มมาจากชุดของการอุทธรณ์ การสมัครสมาชิก การตั้งรหัสผ่านใหม่ การจองห้อง และคุณลักษณะด้านความปลอดภัย \\
ตรวจชนิดข้อมูลฝั่ง frontend & ผ่าน & ไม่มีข้อผิดพลาด \\
ชุดทดสอบส่วนประกอบฝั่ง frontend & ผ่าน 105 จาก 105 & 17 ไฟล์ รายงานครั้งที่ 4 แจ้ง 33 กรณีใน 3 ไฟล์ \\
ชุดตรวจเกณฑ์การยอมรับ & ผ่าน 49 จาก 49 & เท่ารอบก่อน \\
ตรวจรูปแบบโค้ดฝั่ง frontend & 5 ข้อผิดพลาด 8 คำเตือน & รอบก่อน 0 ข้อผิดพลาด ทั้ง 5 ข้ออยู่ในไฟล์ชุดทดสอบการเข้าสู่ระบบที่เพิ่มรอบนี้ และเป็นเรื่องการประกาศชนิดข้อมูลเป็น any ระบบตรวจอัตโนมัติไม่มีขั้นตอน lint จึงไม่จับ \\
ชุดทดสอบผ่านเบราว์เซอร์ & ผ่าน 20 ไม่ผ่าน 5 จาก 25 & รายงานครั้งที่ 4 แจ้ง 13 กรณี รายละเอียดของ 5 กรณีที่ไม่ผ่านอยู่ในตารางถัดไป \\
\bottomrule
\end{tabularx}

กรณีที่ไม่ผ่านทั้ง 5 กรณีไม่ใช่ข้อผิดพลาดของเซิร์ฟเวอร์ แต่เป็นชุดทดสอบที่ค้นหาองค์ประกอบบนหน้าจอด้วยชื่อเดิม ขณะที่หน้าจอถูกแก้โดยตั้งใจในรอบนี้

\noindent\begin{tabularx}{\dimexpr\textwidth-\leftskip\relax}{@{} L{4.4cm} Y @{}}
\toprule
\hrow \thead{กรณีที่ไม่ผ่าน} & \thead{สาเหตุ} \\
\midrule
หน้าผู้ดูแลระบบ ทะเบียนงานตามเวลา & ชุดทดสอบมองหางาน computeAvailability ซึ่งถูกถอดออกโดยตั้งใจในรอบนี้ \\
หน้ารายละเอียดอุปกรณ์ & ชุดทดสอบมองหาหัวข้อของคงเหลือ 14 วัน ซึ่งถูกเอาออกโดยตั้งใจในรอบนี้ \\
หน้ารายการห้อง & ชุดทดสอบมองหาปุ่มจองห้องด้วยชื่อเดิม ขณะที่หน้ารายการห้องถูกออกแบบใหม่ \\
การส่งคำขอยืม 6.1 & ชุดทดสอบมองหาหัวข้อชื่ออุปกรณ์บนร่างคำขอ ขณะที่ร่างคำขอถูกออกแบบใหม่ \\
การส่งคำขอยืม 6.11 & ช่องวันที่รับของเปลี่ยนจากช่องกรอกเป็นปุ่มเลือก ตามงานเลือกอุปกรณ์ตามช่วงเวลา ชุดทดสอบยังพยายามพิมพ์ลงไป \\
\bottomrule
\end{tabularx}

สองกรณีสุดท้ายคือชุดทดสอบของการส่งคำขอยืม ซึ่งเป็นชุดเดียวกับที่เพิ่มเข้ามาในรายงานครั้งที่ 4 เพื่อตอบจุดวัดข้อ 2 ของรอบนั้น ชุดนี้ผ่านตอนที่รวมเข้า main วันที่ 17 กันยายน และเสียไประหว่างรอบนี้เมื่อหน้าจอถูกออกแบบใหม่ ไม่มีใครทราบจนกระทั่งคณะทำงานรันเองในวันรายงาน เพราะระบบตรวจอัตโนมัติไม่เรียกชุดทดสอบผ่านเบราว์เซอร์ ซึ่งเป็นจุดวัดข้อ 3 ที่ไม่ผ่าน

ข้อควรระวังที่คณะทำงานบันทึกไว้คือการ seed ฐานข้อมูลสำหรับชุดทดสอบผ่านเบราว์เซอร์ต้องใช้คำสั่ง seed ที่สร้างบัญชี test\_borrower ไม่ใช่คำสั่ง db:seed ซึ่งเป็นคนละชุดข้อมูลและไม่มีบัญชีนั้น ถ้าใช้ผิดชุด ทุกกรณีจะไม่ผ่านที่ขั้นเข้าสู่ระบบ

% =============================================================================
\clearpage
\section{แผนการดำเนินงานถัดไป}
\label{sec:next}

\subsection{ส่วนที่ยืดออกของ Sprint 4 (25 กันยายน – 1 ตุลาคม 2569)}

สัปดาห์นี้คือส่วนที่ยืดออกของ Sprint 4 ไม่ใช่ Sprint 5 จุดประสงค์คือไล่ตามงานที่ยังไม่เสร็จให้ทันเป้าหมายเดิมที่ 1 ตุลาคม และรับงาน 2 รายการที่ดึงกลับเข้าขอบเขต

\noindent\begin{tabularx}{\dimexpr\textwidth-\leftskip\relax}{@{} L{4.2cm} Y @{}}
\toprule
\hrow \thead{ลำดับ} & \thead{งาน} \\
\midrule
ก่อนอื่น & แก้ชุดทดสอบผ่านเบราว์เซอร์ 5 กรณีที่ไม่ผ่าน ให้ค้นหาองค์ประกอบตรงกับหน้าจอปัจจุบัน \\
ควบคู่กัน & ปิดข้อบกพร่องของหน้าจอที่ยังไม่แก้ 9 ข้อ ตามหัวข้อ \ref{sec:uireview} และกำกับสถานะกลับไปที่เอกสารตรวจทุกข้อที่ปิด \\
ถัดมา & ให้ระบบตรวจอัตโนมัติเรียกชุดทดสอบผ่านเบราว์เซอร์ และขวางการรวมโค้ดเมื่อไม่ผ่าน พร้อมเพิ่มขั้นตอน lint \\
ถัดมา & ตะกร้าข้ามหน่วยงาน ซึ่งยังไม่มีโค้ดเลยทั้งสองฝั่ง และการบังคับเพดานจองล่วงหน้า 90 วันที่ฝั่งเซิร์ฟเวอร์ ซึ่งหน้าเว็บจำกัดไว้แล้วฝ่ายเดียว \\
ท้ายสุด & เพิ่มชุดทดสอบของวงจรยืม–คืนให้เดินครบเส้นทาง และเติมคำแปลไทยของรหัสข้อผิดพลาดที่ขาด 7 ตัว \\
\bottomrule
\end{tabularx}

\subsection{Sprint 5 (2 – 8 ตุลาคม 2569)}

Sprint 5 เป็น Sprint สุดท้ายและไม่รับงานเพิ่มความสามารถตามที่ประกาศไว้ตั้งแต่ต้น งานทั้งหมดในช่วงนี้เป็นการทดสอบถดถอยทั้งระบบ การนำขึ้นใช้งานจริง คู่มือผู้ใช้และคู่มือติดตั้ง และการซ้อมนำเสนอ

งานที่ตัดออกตามแผนสำรองในรายงานครั้งที่ 4 คือตะกร้าข้ามหน่วยงานและการจองล่วงหน้า ถูกดึงกลับเข้าขอบเขตแล้ว เพราะมีเวลาเพิ่มขึ้นหนึ่งสัปดาห์และความสามารถฝั่งเซิร์ฟเวอร์กลับมาเดินตามปกติ ขอบเขตที่ส่งมอบจึงกลับไปเท่ากับข้อเสนอโครงการ

\subsection{จุดวัดถัดไป}

\begin{enumerate}
\item ชุดทดสอบผ่านเบราว์เซอร์ผ่านครบทั้ง 25 กรณี ภายใน 26 กันยายน 2569
\item ระบบตรวจอัตโนมัติเรียกชุดทดสอบผ่านเบราว์เซอร์ และขวางการรวมโค้ดเมื่อไม่ผ่าน ภายใน 27 กันยายน 2569
\item วงจรยืม–คืนเดินครบเส้นทางปกติในชุดทดสอบผ่านเบราว์เซอร์ ตั้งแต่ส่งคำขอถึงตรวจสภาพ ภายใน 29 กันยายน 2569
\item ข้อบกพร่องของหน้าจอที่ยังไม่แก้ 9 ข้อ ปิดครบ และกำกับสถานะในเอกสารตรวจแล้ว ภายใน 29 กันยายน 2569
\item ตะกร้าข้ามหน่วยงานและการจองล่วงหน้าใช้งานได้ทั้งสองฝั่ง ภายใน 1 ตุลาคม 2569
\item ระบบขึ้นใช้งานบนเครื่องจริงและเข้าถึงได้จากภายนอก ภายใน 6 ตุลาคม 2569
\end{enumerate}

\subsection{การประเมินกำหนดเสร็จ}

คณะทำงานประเมินว่าโครงการเสร็จได้ภายใน 8 ตุลาคม 2569 ซึ่งเป็นกำหนดใหม่ และขอบเขตที่ส่งมอบกลับไปเท่ากับข้อเสนอโครงการแล้ว หลังดึงตะกร้าข้ามหน่วยงานและการจองล่วงหน้ากลับเข้ามา

ระยะห่างจากแผน 10 จุดเป็นค่าที่กว้างที่สุดตั้งแต่เริ่มโครงการ แต่ต้องอ่านคู่กับที่มาของมัน ตัวเลขนี้เทียบกับเป้าหมายเดิมที่ไม่ได้เลื่อนตามกำหนดเสร็จ และถูกกดลงจากการหักคุณภาพหน้าจอกับการขยายขอบเขตกลับ ถ้าวัดด้วยวิธีเดียวกับรายงานครั้งที่ 4 ทุกประการ ค่าของวันนี้จะเป็น 90\% ไม่ใช่ \actualpct

เหตุผลที่คณะทำงานยังมั่นใจมีสามข้อ ข้อแรก งานที่เหลือเป็นงานที่ประเมินเวลาได้ คือปิดข้อบกพร่องที่ระบุไว้ครบทุกข้อแล้ว และต่อหน้าจอเข้ากับ procedure ที่มีอยู่ ไม่มีงานที่ยังไม่รู้วิธีทำ ข้อที่สอง ความสามารถฝั่งเซิร์ฟเวอร์กลับมาเดินเต็มที่ในรอบนี้ หลังหยุดนิ่งทั้งรอบก่อน ข้อที่สาม สัปดาห์ที่เพิ่มมาเป็นเวลาไล่งานล้วน ไม่มีเป้าหมายใหม่มาแย่งเวลา

ความเสี่ยงที่เหลืออยู่ข้อเดียวยังเป็นเรื่องเดิม คือระบบตรวจอัตโนมัติยังไม่เรียกชุดทดสอบผ่านเบราว์เซอร์มาสามรอบติดกัน และรอบนี้แสดงผลของการไม่มีกลไกนั้นออกมาแล้ว เมื่อขอบเขตเพิ่งถูกขยายกลับและเหลือเวลาสองสัปดาห์ การแก้ของใหม่จะทำให้ของเก่าเสียโดยไม่มีใครรู้ จุดวัดข้อที่ 2 ของรอบถัดไปจึงตั้งไว้ที่เรื่องนี้โดยตรง

เหตุผลที่ความมั่นใจสูงขึ้นคือความเสี่ยงสองข้อที่รายงานครั้งที่ 4 ระบุไว้ ปิดลงแล้วทั้งคู่ ข้อแรกคือความสามารถใหม่ฝั่ง backend ซึ่งรอบก่อนหยุดนิ่งทั้งรอบ กลับมาเดินและปิดงานที่ค้างนานที่สุดสองเรื่องคือการอุทธรณ์และการจองห้อง ข้อที่สองคือความไม่ตรงกันของสัญญาระหว่างสองฝั่ง ซึ่งเหลือ 0 รายการ

ความเสี่ยงที่เหลืออยู่ข้อเดียวคือการทดสอบ ระบบตรวจอัตโนมัติยังไม่เรียกชุดทดสอบผ่านเบราว์เซอร์มาสามรอบติดกัน และรอบนี้แสดงผลของการไม่มีกลไกนั้นออกมาแล้ว คือชุดทดสอบเสียไปกลางรอบโดยไม่มีใครทราบ เมื่อ Sprint 5 เป็นช่วงทดสอบถดถอยทั้งระบบ การไม่มีกลไกนี้จะทำให้ความถดถอยที่เกิดขึ้นในช่วงสุดท้ายไม่ถูกจับ จุดวัดสองข้อแรกของรอบถัดไปจึงตั้งไว้ที่เรื่องนี้โดยตรง

% =============================================================================
\clearpage
\section*{ภาคผนวก — ค่าที่ใช้พล็อตกราฟ}

ค่าเป้าหมายรายมิติ ณ วันที่กำหนดไว้ คูณด้วยน้ำหนักของมิตินั้นแล้วรวมกัน คอลัมน์ ``รวม'' คือค่าที่พล็อตบนเส้นแผนในรูปที่ \ref{fig:scurve} ตารางนี้เป็นตารางเดียวกับรายงานครั้งที่ 2 ถึงครั้งที่ 4 ทุกช่อง

กำหนดเสร็จที่เลื่อนออกไปหนึ่งสัปดาห์ไม่ได้ทำให้ตารางนี้เปลี่ยน ค่าเป้าหมายยังผูกกับวันเดิม และยังแตะ 100\% ที่ 1 ตุลาคม ช่วง 2 ถึง 8 ตุลาคมไม่มีแถวใหม่ เพราะไม่มีเป้าหมายใหม่ในช่วงนั้น มีแต่เวลาให้ไล่ตามงานที่ยังไม่เสร็จ วันที่กำกับหัวแถวจึงเป็นวันของเป้าหมาย ไม่ใช่ขอบ Sprint ซึ่งขยับไปแล้วตามหัวข้อ \ref{sec:overview}

\noindent\begin{tabularx}{\dimexpr\textwidth-\leftskip\relax}{@{} L{3.2cm} L{1.0cm} L{1.0cm} L{1.0cm} L{1.0cm} L{1.0cm} L{1.0cm} Y @{}}
\toprule
\hrow \thead{ขอบ Sprint} & \thead{BE Des} & \thead{FE Des} & \thead{BE Dev} & \thead{FE Dev} & \thead{Test} & \thead{Doc} & \thead{รวม} \\
\hrow \thead{} & \thead{10\%} & \thead{10\%} & \thead{30\%} & \thead{25\%} & \thead{15\%} & \thead{10\%} & \thead{100\%} \\
\midrule
เริ่ม 24 ก.ค. & 0 & 0 & 0 & 0 & 0 & 10 & 1 \\
S0 · 30 ก.ค. & 45 & 30 & 5 & 5 & 0 & 20 & 12 \\
S1 · 13 ส.ค. & 88 & 75 & 25 & 42 & 10 & 65 & 42 \\
S2 · 27 ส.ค. & 95 & 85 & 60 & 65 & 35 & 70 & 65 \\
S3 · 10 ก.ย. & 100 & 92 & 78 & 78 & 52 & 76 & 78 \\
$*$ 24 ก.ย. & 100 & 100 & 97 & 97 & 80 & 88 & 94 \\
1 ต.ค. & 100 & 100 & 100 & 100 & 100 & 100 & 100 \\
8 ต.ค. (กำหนดใหม่) & 100 & 100 & 100 & 100 & 100 & 100 & 100 \\
\bottomrule
\end{tabularx}

แถวที่มีเครื่องหมายดอกจันคือวันรายงานรอบนี้ ซึ่งตรงกับวันของเป้าหมายพอดี จึงอ่านค่าได้ตรง ไม่ต้องสอดแทรกเหมือนรายงานครั้งที่ 2 และครั้งที่ 4

ค่าผลจริง ณ วันรายงาน คือ 100, 92, 88, 82, 62 และ 88 ตามลำดับมิติเดียวกัน ถ่วงน้ำหนักแล้วได้ 84.20 ปัดเป็น \actualpct

ถ้าไม่หักข้อบกพร่องของหน้าจอตามหัวข้อ \ref{sec:uireview} ค่ารายมิติจะเป็น 100, 100, 90, 90, 66 และ 88 ถ่วงน้ำหนักแล้วได้ 88.20 ปัดเป็น 88\%

ถ้าไม่หักข้อบกพร่อง และยังไม่ดึงงาน 2 รายการกลับเข้าขอบเขตด้วย ค่าจะเป็น 100, 100, 93, 92, 66 และ 88 ถ่วงน้ำหนักได้ 89.60 ปัดเป็น 90\% ซึ่งเป็นค่าที่วิธีนับของสี่ฉบับก่อนบนขอบเขตเดิมจะให้

ค่าผลจริงรอบก่อนหน้าที่อยู่บนเส้นผลจริง คือ 52\% ณ 27 สิงหาคม 70\% ณ 4 กันยายน 75\% ณ 10 กันยายน และ 79\% ณ 17 กันยายน ตามที่ส่งในรายงานครั้งที่ 1 ถึงครั้งที่ 4 ไม่มีการแก้ย้อนหลัง

\end{document}
